\documentclass[11pt,a4paper]{article}
\usepackage[T1]{fontenc}
\usepackage[utf8]{inputenc}
\usepackage{lmodern}
\usepackage[margin=2.6cm]{geometry}
\usepackage{amsmath,amssymb,amsthm}
\usepackage{booktabs}
\usepackage{array}
\usepackage[hidelinks]{hyperref}

\newtheorem{theorem}{Theorem}[section]
\newtheorem{lemma}[theorem]{Lemma}
\newtheorem{proposition}[theorem]{Proposition}
\newtheorem{corollary}[theorem]{Corollary}
\theoremstyle{definition}
\newtheorem{definition}[theorem]{Definition}
\theoremstyle{remark}
\newtheorem{remark}[theorem]{Remark}

\newcommand{\Q}{\mathbb{Q}}
\newcommand{\E}{\mathbb{E}}
\newcommand{\PP}{\mathbb{P}}
\newcommand{\cR}{\mathcal{R}}
\newcommand{\cI}{\mathcal{I}}
\newcommand{\sep}{\operatorname{sep}}
\newcommand{\sepw}{\operatorname{sep}_{\theta}}
\newcommand{\dfc}{\operatorname{def}}
\newcommand{\tw}{\operatorname{tw}}
\newcommand{\rank}{\operatorname{rank}}
\newcommand{\Gr}{\Gamma}
\newcommand{\nocc}{n_{\mathrm{occ}}}
\newcommand{\Nbar}{\overline{N}}
\newcommand{\DL}{DL}

\title{Measurement Rows Absorb at Most a \texorpdfstring{$\sqrt{s}$}{sqrt(s)} Factor of the Cut:\\
A Determinant Bound for Clause-Product Interpolation Schemes}
\author{Lucas Requena\\[2pt]
{\small Computer Science student}\\
{\small \texttt{l.requena@outlook.fr}}}
\date{November 2026}

\begin{document}
\maketitle

\begin{abstract}
Dias and Lima (arXiv:2611.00517) abstract the clause-product form of the interpolation
algorithm in OpenAI's claimed refutation of the Exponential Time Hypothesis (ETH) as a
\emph{scheme}: at each evaluation point the variables are partitioned into groups of size at
most $s$, and a clause may see the other groups only through linear \emph{measurement rows}.
They prove a cost of $2^{\Omega(\sqrt n)}$ on expanding formulas and leave open, in their
Section~9, whether measurements can ``absorb the cut''. We show that measurements absorb at
most a $\sqrt{s}$ factor of it: at some point, the number of rows plus the Taylor order
exceeds a constant times $\sep(F)/\sqrt{s}$, where $\sep(F)$ is the balanced-separator number
of the primal graph. For deformations that depend only on the variable, the factor
$\sqrt s$ disappears. The proof uses an $n\times n$ determinant built from one occurrence
per variable, and an anticoncentration bound. Consequently, on bounded-occurrence formulas
with linear separators and on random 3-CNF of density above $3$, every scheme applied to the
formula itself costs $2^{\Omega(n^{2/3})}$ in the product-type cost model. We also observe
that the DP-type cost model is degenerate: a trivial scheme has charge
$2^{\lceil\sqrt n\rceil}$ on every formula. The deterministic part of the proof is
machine-checked in Lean~4. The results were obtained with AI tools, have not yet been
reviewed by human experts, and say nothing about the preprint's sharper routes, about ETH,
or about P versus NP.
\end{abstract}

\section{Introduction}\label{sec:intro}

The Exponential Time Hypothesis (ETH) of Impagliazzo and Paturi~\cite{IP01} states that 3-SAT
on $n$ variables needs time $2^{\delta n}$ for some $\delta>0$. In October 2026 an OpenAI
preprint~\cite{OpenAI} claimed a deterministic algorithm that counts the satisfying assignments
of a $k$-CNF formula in time $2^{O_k(n\log\log n/\sqrt{\log n})}$, which would refute ETH. An
exposition was written by O'Donnell~\cite{ODonnell}, and a second preprint, by Anthropic's
Claude~\cite{Claude}, announced a $2^{o(n)}$-time algorithm for 3-SAT. We do not assess the
correctness of either preprint, and we do not study the second one.

Dias and Lima~\cite{DL} analysed the OpenAI preprint. Among other things, they isolate the
\emph{clause-product form} of its interpolation algorithm and abstract it as a
\emph{clause-product interpolation scheme} (their Definition~5.1, reproduced in
Section~\ref{sec:setting}). In a cost model that charges what the preprint's evaluation pays,
they prove that every such scheme costs $2^{\Omega(\sqrt{\tw})}$ on formulas of primal
treewidth $\tw$, hence $2^{\Omega(\sqrt n)}$ on expanding formulas. The preprint's own
clause-product design costs $2^{O(n(\log\log n)^2/\sqrt{\log n})}$ on bounded-occurrence
formulas, so a gap remains. In their Section~9, Dias and Lima explain that their technique
cannot go beyond $2^{\tilde O(\sqrt n)}$, and they state what is needed to narrow the gap from
below: ``one must show that measurements cannot absorb the cut''. They make a pairwise
observation for groups with one measurement row (a $2\times2$ minor) and write that ``turning
this into a bound for groups with several measurement rows is open''.

\paragraph{Our result.} We give such a bound. Let $F$ be a CNF formula on $n$ variables in
which every variable occurs at most $b$ times, and let $\sep(F)$ be the minimum size of a set
of variables whose removal leaves components of at most $2n/3$ vertices in the primal graph.
In the notation of Definition~5.1 of~\cite{DL}, let $s_t$ be the largest group size at the
point $t$, let $\rho_t$ be the total number of independent measurement rows, and let $q_t$ be
the Taylor order. Theorem~\ref{thm:general} states that every scheme for $F$ has a point $t$
with
\begin{equation}\label{eq:intro}
\rho_t+q_t+1\;>\;\min\Bigl(c_b\,\frac{\sep(F)}{\sqrt{s_t}},\ \frac{2n}{3}-s_t\Bigr),
\qquad c_b=\frac{(1-e^{-2})(b-1)}{4\sqrt2\,b^2},\quad c_3=\frac{1-e^{-2}}{18\sqrt2}=0.03396\ldots
\end{equation}
If the deformation of an occurrence depends only on its variable and not on its clause,
Theorem~\ref{thm:uniform} gives the same statement with $\sep(F)$ in place of
$c_b\sep(F)/\sqrt{s_t}$. So when $\sep(F)=\Omega(n)$, a scheme needs about one measurement row
(or one unit of Taylor order) per $\sqrt{s_t}$ cut variables in general, and one per cut
variable in the variable-uniform case. We do not know whether the factor $\sqrt{s_t}$ can
really be gained by a clause-dependent design; this is open.

\paragraph{Contributions.}
\begin{enumerate}
\item \textbf{A witness-determinant lemma} (Lemma~\ref{lem:det}). Choose one occurrence of
 each variable and form the $n\times n$ polynomial matrix whose rows are the chosen
 deformation vectors. Its determinant is a nonzero polynomial of degree at most $nR$, and its
 order of vanishing at an evaluation point is at least a combinatorial \emph{deficiency} of
 the design at that point minus the number of measurement rows. The $2\times2$ minor of
 \cite[Section~9]{DL} is a $2\times 2$ minor of this matrix; the full determinant handles any
 number of rows per group.
\item \textbf{Variable-uniform deformations gain nothing from measurements}
 (Theorem~\ref{thm:uniform}): some point has $\rho_t+q_t+1>\min(\sep(F),2n/3-s_t)$.
\item \textbf{General deformations gain at most a $\sqrt{s}$ factor}
 (Theorem~\ref{thm:general}), by averaging the determinant lemma over random choices of
 occurrences and bounding the expected deficiency from below with an anticoncentration
 inequality with explicit constant (Lemmas~\ref{lem:anti} and~\ref{lem:expdef}).
\item \textbf{Consequences} (Corollary~\ref{cor:main}) for bounded-occurrence formulas with
 $\sep(F)\ge cn$ and for random 3-CNF with $m=\Delta n$ clauses, $\Delta>3$: the design
 described in \cite[Section~9]{DL} (group size, Taylor order and number of rows all
 $\tilde O(\sqrt n)$) does not exist, and in the product-type cost model of
 \cite[Section~5.2]{DL} every scheme applied to the formula itself costs
 $2^{\Omega(n^{2/3})}$; variable-uniform schemes cost $2^{\Omega(n)}$.
\item \textbf{The DP-type cost model is degenerate} (Proposition~\ref{prop:dp}). In the second
 cost model of \cite[Section~5.2]{DL}, a trivial scheme (constant deformation, fully measured
 groups of size $\lceil\sqrt n\rceil$) has charge $2^{\lceil\sqrt n\rceil}$ on every formula.
 So the bound $2^{\Omega(\sqrt n)}$ of~\cite{DL} is tight in that model for a trivial reason,
 no technique can prove more in it, and the question of Section~9 is meaningful only in a
 model that charges rows, such as the product-type model.
\item \textbf{Machine verification} (Section~\ref{sec:lean}). The determinant lemma, the
 Hermite counting step, the averaging step, the separator lemma, Theorem~\ref{thm:uniform},
 and the deduction of Theorem~\ref{thm:general} from the expected-deficiency lemma are
 formalised in Lean~4 with Mathlib. The expected-deficiency lemma itself, the lemma on random
 3-CNF and the corollary are proved on paper only.
\item \textbf{Computational checks} (Section~\ref{sec:checks}) on small instances, which can
 detect errors but prove nothing asymptotic.
\end{enumerate}

\paragraph{What this paper does not show.}
\begin{itemize}
\item The cost statements are confined to the \emph{product-type} cost model
 of~\cite[Section~5.2]{DL}. They are lower bounds on the charges of that model, and they are
 statements about running time only to the extent that these charges are paid. In the
 DP-type model nothing beyond $2^{\lceil\sqrt n\rceil}$ can be proved
 (Proposition~\ref{prop:dp}).
\item The results concern a scheme applied to $F$ \emph{itself}. The preprint applies its
 scheme to a bounded-occurrence formula $F^\star$ obtained from $F$ by variable copies, and
 to the leaves of a sparsification. Theorem~8.3 of~\cite{DL} covers these; our
 Corollary~\ref{cor:main} does not, because we have not proved $\sep(F^\star)=\Omega(n^\star)$
 for those formulas.
\item The results say nothing about the preprint's parity form or its tensor contraction,
 which are outside Definition~5.1 (see \cite[Section~5.1]{DL}), nothing about the algorithm
 of~\cite{Claude}, and nothing about the truth of ETH.
\item The results have no bearing on P versus NP. They are barriers for one abstraction of
 one algorithmic route, not for a problem.
\item We give no upper bound. Whether some design attains $\rho_t\approx n/\sqrt{s_t}$, and
 hence whether $n^{2/3}$ is the right exponent in the product-type model, is open.
\item No human expert has reviewed this work yet (Section~\ref{sec:ai}).
\end{itemize}

\paragraph{The idea in plain terms.}
At an evaluation point the variables are split into groups. A clause assigned to one group
may read the variables of another group only through that group's measurement rows. Pick, for
every variable $v$, one clause containing it, and look at the linear form that this clause
reads in place of $x_v$. At $T=0$ this form is $x_v$ itself, so the $n$ forms are
independent and their determinant is $1$. At an evaluation point $t$, consider one group $S$
with $a$ variables. The forms attached to clauses of other groups can only contribute, on the
coordinates of $S$, vectors from the row space of $S$. So if fewer than $a$ of the chosen
clauses are assigned to $S$, and $S$ has few rows, the $n$ forms become dependent at $t$: the
determinant vanishes there, to an order that counts the missing dimensions. The determinant
has degree at most $MR$, where $M$ is the number of occurrences, so its orders of vanishing
at the evaluation points sum to at most $MR$. On the other hand, a recovery rule that is
exact for all polynomials of degree $MR$ must use more than $MR$ pieces of Hermite data in
total. Hence at some point the missing dimensions, the \emph{deficiency}, must be paid for
by rows or by Taylor order.

How large is the deficiency? If the deformation depends only on the variable, every clause
containing $v$ reads the same form, and a variable whose clauses are assigned to two groups
is constrained everywhere; the deficiency is then the number of such \emph{split} variables,
and these form a separator of the primal graph. If the deformation depends on the clause, a
design may hope that the chosen clauses happen to fill every group exactly. But the choice of
clauses is ours: choosing them at random, the number of chosen clauses assigned to a group
fluctuates by about the square root of its variance, and these fluctuations add up to at
least (number of split variables)$/\sqrt{s}$. This is where the $\sqrt s$ comes from. It is a
loss of the proof technique, and we do not know whether it reflects a real possibility.

\paragraph{Organisation.} Section~\ref{sec:related} discusses related work and compares our
bounds with those of~\cite{DL}. Section~\ref{sec:setting} reproduces the definitions
of~\cite{DL} and fixes notation. Section~\ref{sec:results} states the results and
Section~\ref{sec:proofs} proves them. Section~\ref{sec:dp} treats the DP-type model.
Sections~\ref{sec:lean} and~\ref{sec:checks} describe the Lean formalisation and the
computational checks. Section~\ref{sec:ai} describes how the work was done, and
Section~\ref{sec:concl} concludes with open problems. Appendix~\ref{app:random} reproduces,
for completeness, the proofs of two facts about random 3-CNF taken from~\cite{DL}.

\section{Related work}\label{sec:related}

\paragraph{ETH and the 2026 preprints.} ETH was formulated by Impagliazzo and
Paturi~\cite{IP01}, and the sparsification lemma of Impagliazzo, Paturi and Zane~\cite{IPZ01}
transfers it to many problems. The OpenAI preprint~\cite{OpenAI} claims subexponential
deterministic counting for $k$-CNF. According to~\cite{DL}, its headline bound is proved by a
contraction of Boolean tensor networks, and its body develops an interpolation algorithm in
two forms, a clause-product form and a parity form. O'Donnell~\cite{ODonnell} wrote an
exposition of the preprint. A separate preprint by Anthropic's Claude~\cite{Claude} announces a
$2^{o(n)}$-time algorithm for 3-SAT. Our results concern neither the tensor contraction, nor
the parity form, nor the algorithm of~\cite{Claude}.

\paragraph{Dias and Lima.} Our starting point is~\cite{DL}. All our definitions of schemes
and cost models are theirs. We use or compare with the following results of~\cite{DL}.
\begin{itemize}
\item \emph{Lemma~6.2 (vanishing product).} The product $\prod_o c_o(T)$ of the own
 coefficients has degree at most $MR$ and equals $1$ at $T=0$, so at some point $t$ at most
 $q_t$ occurrences lose their own variable. Our Lemma~\ref{lem:hermite} uses the same
 argument, applied to the polynomial $\prod_{t}(T-t)^{q_t+1}$; the argument is theirs. Our
 Lemma~\ref{lem:det} replaces the product of the $M$ diagonal entries $c_o$ by the
 determinant of an $n\times n$ matrix, of degree at most $nR$.
\item \emph{Theorem~7.1(ii).} At the point of their Lemma~6.2, deleting at most
 $q_t+b\,\mu_ts_t$ edges of the incidence graph leaves every component inside one group,
 where $\mu_t$ is the number of measured groups. If $s_t\le 2n/3$, the variables incident to
 the deleted edges form a $2/3$-separator of the primal graph\footnote{If $u,v$ are adjacent
 in the primal graph and neither is incident to a deleted edge, they share a clause $e$ and
 both edges $(e,u),(e,v)$ survive, so $u,v$ lie in the same component after the deletion,
 hence in the same group.}, so this gives
 \[
 q_t+b\,\mu_t s_t\ \ge\ \sep(F).
 \]
 This bound charges \emph{measured groups}: a measured group of size $s_t$ may absorb all of
 its roughly $bs_t$ foreign occurrences, whatever its number of rows. Since
 $\rho_t\ge\mu_t$, it implies $\rho_t\ge(\sep(F)-q_t)/(bs_t)$. Our
 Theorem~\ref{thm:general} charges \emph{rows} and gives
 $\rho_t>c_b\sep(F)/\sqrt{s_t}-q_t-1$ when $c_b\sep(F)/\sqrt{s_t}\le 2n/3-s_t$. In short:
 $\mu_ts_t\gtrsim n$ in~\cite{DL}, against $\rho_t\sqrt{s_t}\gtrsim n$ here, on formulas with
 $\sep(F)=\Omega(n)$ and at possibly different points $t$.
\item \emph{Theorem~8.1 (treewidth barrier) and Lemma~8.2.} In the cost model, every DP-type,
 hence every product-type, scheme on $F$ costs at least
 $2^{\min(\sqrt{(\tw+1)/2}-1,\,(\tw+1)/2,\,\lfloor n/2\rfloor)}$ with $\tw$ the primal
 treewidth, and random 3-CNF with $\Delta>3$ has linear primal treewidth w.h.p. A graph of
 treewidth $\tw$ has a separator of at most $\tw+1$ vertices leaving components of at most
 $n/2$ vertices~\cite{RS86,Bod93}, so $\sep(F)\le\tw+1$, and Theorem~8.1 of~\cite{DL}
 applies to the families of our Corollary~\ref{cor:main}. It gives $2^{\Omega(\sqrt n)}$ for
 a scheme applied to $F$ itself. This is the fair comparison for our
 $2^{\Omega(n^{2/3})}$, which holds in the product-type model only.
 Our Lemma~\ref{lem:random} adapts the proof of their Lemma~8.2 from treewidth to
 $2/3$-separators and uses the two facts of their Appendix~A.
\item \emph{Theorem~8.3 (random 3-CNF).} W.h.p. the clause-product scheme costs
 $2^{\Omega_\Delta(\sqrt n)}$ when applied (a) to the preprint's bounded-occurrence formula
 $F^\star$ and (b) to the leaves of its sparsification. Our results do \emph{not} improve
 this theorem: we treat $F$ itself, and we have not established the separator hypothesis for
 $F^\star$ or for the leaves.
\end{itemize}
Table~\ref{tab:compare} in Section~\ref{sec:results} summarises the comparison.

\paragraph{Tools.} The fact that the determinant of a polynomial matrix vanishes at a point
to order at least the dimension of the kernel there is classical; it follows from the local
Smith form~\cite{GLR82} and is a matrix analogue of the multiplicity lemmas of the polynomial
method~\cite{DKSS13}. We include the short proof. Anticoncentration of sums of independent
random variables goes back to Littlewood and Offord~\cite{LO43}, Erd\H{o}s~\cite{Erd45},
Kolmogorov~\cite{Kol58} and Rogozin~\cite{Rog61}. The lower bound
$\E|Z|\ge(\E Z^2)^{3/2}/(\E Z^4)^{1/2}$ that we use is the fourth-moment method, which goes
back to Littlewood~\cite{Lit30}; see Berger~\cite{Ber97}. Balanced separators were introduced
by Lipton and Tarjan~\cite{LT79}; their relation to treewidth is due to Robertson and
Seymour~\cite{RS86}, see also~\cite{Bod93}. Expansion of random CNF formulas is the standard
source of lower bounds in proof complexity~\cite{CS88,BSW01,BSI10}; as in~\cite{DL}, we use
only two elementary properties of random 3-CNF, proved in Appendix~\ref{app:random}.

\paragraph{Barriers for methods.} Like those of~\cite{DL}, our results are barriers for a
method in a restricted model, in the spirit of the limits of the laser method for matrix
multiplication~\cite{AFL15,AVW18}. They are not lower bounds for a problem.

\paragraph{Novelty.} We are not aware of a prior proof of
Theorems~\ref{thm:uniform} and~\ref{thm:general}. This assessment rests on searches made
with AI tools, not on a systematic survey, and the source~\cite{DL} is itself a few weeks
old. We welcome corrections.

\section{Setting and notation}\label{sec:setting}

Sections~\ref{sec:def}--\ref{sec:question} reproduce material from~\cite{DL}; indented
passages are quotations. Section~\ref{sec:ournot} introduces our own notation.

\subsection{Clause-product interpolation schemes}\label{sec:def}

Each clause $e$ has a satisfaction polynomial $\varphi_e$, equal to $1$ when the clause is
satisfied and $0$ otherwise, so that
$\#F=\sum_{x\in\{0,1\}^n}\prod_e\varphi_e(x)$ \cite[(1)]{DL}. The deformed sum is
\begin{equation}\tag{DL.2}\label{eq:H}
H(T)=\sum_{x\in\{0,1\}^n}\ \prod_e\varphi_e\bigl(\langle P_{e,v}(T),x\rangle\bigr)_{v\in e},
\end{equation}
which is equation~(2) of~\cite{DL}. The following is quoted from \cite[Section~5]{DL}, with
their equation~(3) renumbered~\eqref{eq:three}.

\begin{quote}
Let $F$ be a CNF over $x_1,\dots,x_n$ with clauses of at most $k$ literals, $I_e$ the variable
set of clause $e$, and $O=\{(e,v):v\in I_e\}$ its occurrences, $M=|O|$. Let $b$ bound the
number of occurrences of a variable. The incidence graph $G_F$ joins each clause to its
variables; the primal graph $\mathrm{Pr}(F)$ joins two variables that share a clause.

\textbf{Definition 5.1} (Clause-product interpolation scheme). A scheme for $F$ consists of
the following.
\begin{itemize}
\item[(a)] \emph{Deformation.} For each occurrence $o=(e,v)$ a vector $P_o(T)\in\Q[T]^n$ of
 degree at most $R$ with $P_o(0)=e_v$; $H(T)$ as in (2). Its own coefficient is
 $c_o(T)=(P_o(T))_v$, a polynomial of degree at most $R$ with $c_o(0)=1$. Since
 $\deg\varphi_e=|I_e|$, $\deg H\le MR$.
\item[(b)] \emph{Recovery.} A finite set $X\subset\Q\setminus\{0\}$, orders $q_t\ge0$ for
 $t\in X$, and a linear functional $\Lambda$ on Hermite data
 $(f^{(j)}(t))_{t\in X,\,j\le q_t}$ with $\Lambda(f)=f(0)$ for every $f\in\Q[T]$ of degree at
 most $MR$.
\item[(c)] \emph{Designs.} For each $t\in X$: a partition of $[n]$ into groups
 $S_{t,1},S_{t,2},\dots$; a group $\alpha_t(e)$ for each clause; a subspace
 $\cR_{t,h}\subseteq\Q^{S_{t,h}}$ (the span of the measurement rows) for each group, such
 that for every clause $e$, every $v\in I_e$ and every group $h\ne\alpha_t(e)$,
 \begin{equation}\tag{DL.3}\label{eq:three}
 P_{e,v}(t)\big|_{S_{t,h}}\in\cR_{t,h}.
 \end{equation}
 Write $s_t=\max_h|S_{t,h}|$, $\rho_t=\sum_h\dim\cR_{t,h}$, and $\mu_t$ for the number of
 measured groups ($\cR_{t,h}\neq0$). The interaction graph $\cI_t$ has the groups as
 vertices; for each clause $e$ the groups
 $\{\alpha_t(e)\}\cup\{h\ne\alpha_t(e):P_{e,v}(t)|_{S_{t,h}}\ne0\text{ for some }v\in I_e\}$
 form a clique. Let $w_t=\tw(\cI_t)$.
\end{itemize}
The clause-product form of the preprint's interpolation algorithm (its Sections 2.1--2.8) is
an instance, with $\cR_{t,h}$ the row space of its integer measurement matrix $C_{t,h}$ and
$\Lambda$ Hermite interpolation with $q_t=\lceil 3nR/|X|\rceil$ (Section 2.5 of [the
preprint]). We say that $o=(e,v)$ loses its own variable at $t$ if $c_o(t)=0$.
\end{quote}

Here $e_v$ is the $v$-th unit vector of $\Q^n$, and $P_{e,v}(t)|_{S_{t,h}}$ is the restriction
of the vector $P_{e,v}(t)\in\Q^n$ to the coordinates in $S_{t,h}$. We always assume $n\ge1$.
Our results use only the following parts of the definition: the vectors $P_o$ with
$\deg P_o\le R$ and $P_o(0)=e_v$; the exactness of $\Lambda$ on all polynomials of degree at
most $MR$; and condition~\eqref{eq:three}. The polynomial $H$, the signs of the literals and
the clause width play no role in the proofs.

\subsection{Cost models}\label{sec:cost}

The following is quoted from \cite[Section~5.2]{DL}.

\begin{quote}
We charge each point $t$ at least
\[
\kappa(t)\ \ge\ \max\Bigl(2^{s_t},\ 2^{w_t},\ \tbinom{n}{\min(q_t,\lfloor n/2\rfloor)}\Bigr)
\qquad\text{(DP-type evaluation)},
\]
and, for the preprint's product-type evaluation, additionally $\kappa(t)\ge2^{\rho_t}$. A
scheme runs in time at most $T$ in this model if $\max_{t\in X}\kappa(t)\le T$. All our lower
bounds hold in this model; they are statements about running time only to the extent that the
model's charges are paid.
\end{quote}

Dias and Lima then justify the charges. For the preprint's evaluation, which enumerates each
group's assignments, every tuple of measurement readings and the Taylor data, the charges are
paid. They show $w_t\le\mu_t\le\rho_t$, so ``the product-type model implies the DP-type
one''. About the DP-type model they write:

\begin{quote}
The DP-type model is meant for an evaluation that still counts inside each group by
enumeration, as the preprint does, but sums over measurement readings by dynamic programming
instead of enumerating all tuples. [\dots] Any elimination order over the readings creates a
table indexed by at least $\tw(G_{\mathcal M})+1\ge w_t$ readings, and a nonzero measurement
row takes at least two values on Boolean inputs, so it costs at least $2^{w_t}$. An evaluator
that also counted inside groups by dynamic programming would not pay $2^{s_t}$, and our model
does not cover it.
\end{quote}

Here $G_{\mathcal M}$ is a graph on the measured groups defined in~\cite{DL} (the torso of
$\cI_t$ on the measured groups); we do not need it.

We call \emph{DP-type charge} of a scheme the number
$\max_{t\in X}\max\bigl(2^{s_t},2^{w_t},\binom{n}{\min(q_t,\lfloor n/2\rfloor)}\bigr)$, and
\emph{product-type charge} the maximum of the DP-type charge and $\max_{t\in X}2^{\rho_t}$.
``A scheme costs at least $C$ in the product-type model'' means that its product-type charge
is at least $C$. All the caveats of~\cite{DL} apply to our statements: they bound the charges
of the model, and an evaluator that does not pay $2^{s_t}$ or $2^{\rho_t}$ is not covered.

\subsection{The question of Section 9 of Dias and Lima}\label{sec:question}

The following is quoted from \cite[Section~9, ``The gap'']{DL}; there $\rho(n)$ denotes the
preprint's exponent $n(\log\log n)^2/\sqrt{\log n}$, which is unrelated to $\rho_t$.

\begin{quote}
On expanding inputs of bounded density, the clause-product scheme is confined between
$2^{\Omega(\sqrt n)}$ (Theorem 8.3) and the preprint's $2^{O(\rho(n))}$, in the cost model.
Our lower-bound technique cannot exceed $2^{\tilde O(\sqrt n)}$: its only constraint is
$(w+1)s+q\gtrsim\tw$, which $s=w=q=\sqrt n$ satisfies. Narrowing the gap needs one of two new
ingredients. From above, a design on an expander with $\sqrt n$ groups, $O(\sqrt n)$
interaction width, Taylor order $O(\sqrt n)$ and $O(1)$ or polylogarithmically many rows per
measured group would bring the scheme to $2^{\tilde O(\sqrt n)}$. From below, one must show
that measurements cannot absorb the cut. A measured group of dimension one can absorb two
foreign occurrences $(e,v),(e',v')$ with $v\ne v'$ at $t$ only if the $2\times2$ minor
$c_{e,v}(t)c_{e',v'}(t)-p_{e,v,v'}(t)p_{e',v',v}(t)$ vanishes; that polynomial equals $1$ at
$T=0$, so each such pair is absorbed at no more than $2R$ points. (Occurrences of the same
variable can be absorbed together at no cost.) Turning this into a bound for groups with
several measurement rows is open.
\end{quote}

The symbol $p_{e,v,v'}$ is not defined in~\cite{DL}; we read it as the coordinate
$(P_{e,v}(T))_{v'}$. Three remarks on our reading of this passage.

\emph{(1) It states a direction, not a formal question.} We interpret the target ``from
above'' as a scheme with $s_t$, $q_t$ and $\rho_t$ all at most $\sqrt n\log^{\ell}n$ at every
point, for some fixed $\ell$: with $\sqrt n$ groups and polylogarithmically many rows per
group, $\rho_t=\tilde O(\sqrt n)$, and the cost $2^{\tilde O(\sqrt n)}$ requires
$s_t=\tilde O(\sqrt n)$ in both cost models because $2^{s_t}$ is charged. This last point
depends entirely on the charge $2^{s_t}$, hence on the caveat quoted in
Section~\ref{sec:cost}.

\emph{(2) The $2\times2$ minor is a pairwise observation.} If $v,v'\in S_{t,h}$,
$\dim\cR_{t,h}=1$ and both clauses $e,e'$ are assigned to other groups, then
\eqref{eq:three} puts the restrictions of $P_{e,v}(t)$ and $P_{e',v'}(t)$ to $S_{t,h}$ on
one line, so the $2\times2$ matrix of their coordinates $v,v'$ is singular. In~\cite{DL}
this is not turned into a bound, even for one row per group.

\emph{(3) The parenthetical sentence explains our construction.} Occurrences of the same
variable can be absorbed together, so only one occurrence per variable can be used in a
determinant. This is why we work with \emph{witness maps}, which select one occurrence of
each variable. The $2\times2$ matrix of (2) is the minor of our witness matrix on rows and
columns $v,v'$, for a witness map that selects $(e,v)$ and $(e',v')$.

\subsection{Our notation}\label{sec:ournot}

Table~\ref{tab:notation} lists the notation. We write $\Gr(F)$ for the primal graph (written
$\mathrm{Pr}(F)$ in~\cite{DL}), $\PP$ for probability and $\E$ for expectation, and
$(y)_+=\max(y,0)$. For a variable $v$, let $O_v=\{(e,v):v\in I_e\}$ be its set of occurrences
and $d_v=|O_v|$. So $M=\sum_vd_v$. Let $\nocc$ be the number of variables with $d_v\ge1$;
clearly $\nocc\le\min(n,M)$.

Fix a scheme and a point $t\in X$. When $t$ is fixed we write $S_h=S_{t,h}$, $a_h=|S_h|$ and
$r_h=\dim\cR_{t,h}$, so that $\sum_ha_h=n$, $\sum_hr_h=\rho_t$ and $a_h\le s_t$.

\begin{definition}[Witness maps, witness matrix]\label{def:witness}
A \emph{witness map} $\omega$ assigns to every variable $v$ with $d_v\ge1$ one of its
occurrences $\omega(v)\in O_v$. The \emph{witness matrix} $\Phi^\omega(T)$ is the $n\times n$
matrix over $\Q[T]$ whose row $v$ is $P_{\omega(v)}(T)$ if $d_v\ge1$, and the constant vector
$e_v$ if $d_v=0$.
\end{definition}

\begin{definition}[Home groups, deficiency]\label{def:def}
Let $t\in X$ and let $\omega$ be a witness map. The \emph{home group} $h^\omega_t(v)$ of a
variable $v$ is $\alpha_t(e)$ if $\omega(v)=(e,v)$; if $d_v=0$, it is the group that
contains $v$. Put $I^\omega_{t,h}=\{v:h^\omega_t(v)=h\}$ and $N^\omega_{t,h}=|I^\omega_{t,h}|$.
The \emph{deficiency} at $t$ is
\[
\dfc^\omega_t=\sum_h\bigl(a_h-N^\omega_{t,h}\bigr)_+ .
\]
\end{definition}

Every variable has exactly one home group, so $\sum_hN^\omega_{t,h}=n$.

\begin{definition}[Split variables]\label{def:split}
A variable $v$ is \emph{split at $t$} if it occurs in two clauses $e,e'$ with
$\alpha_t(e)\ne\alpha_t(e')$. Let $\Sigma_t$ be the set of split variables; every
$v\in\Sigma_t$ has $d_v\ge2$. A variable $v\notin\Sigma_t$ has a home group $h_t(v)$ that
does not depend on $\omega$: if $d_v\ge1$ it is the common value of $\alpha_t(e)$ over the
clauses $e\ni v$, and if $d_v=0$ it is the group that contains $v$. Put
$J_{t,h}=\{v\notin\Sigma_t:h_t(v)=h\}$. The sets $J_{t,h}$ partition $[n]\setminus\Sigma_t$.
\end{definition}

\begin{definition}[Separators]\label{def:sep}
A set $Y\subseteq[n]$ is a \emph{$2/3$-separator} of $F$ if every connected component of
$\Gr(F)-Y$ has at most $2n/3$ vertices. Let $\sep(F)$ be the minimum size of a
$2/3$-separator. Let $\theta(d)=2(d-1)/d^2$, and
\[
\sepw(F)=\min\Bigl\{\sum_{v\in Y}\theta(d_v)\ :\ Y\text{ a }2/3\text{-separator of }F
\text{ with }d_v\ge2\text{ for all }v\in Y\Bigr\},
\]
with $\sepw(F)=+\infty$ if there is no such $Y$, and the convention $\min(+\infty,y)=y$.
\end{definition}

The function $\theta$ is decreasing on $d\ge2$, because
$\frac{d}{dd}\frac{d-1}{d^2}=\frac{2-d}{d^3}\le0$. Hence, if $d_v\le b$ for all $v$ with
$b\ge2$,
\begin{equation}\label{eq:sepw}
\sepw(F)\ \ge\ \theta(b)\,\sep(F)=\frac{2(b-1)}{b^2}\,\sep(F).
\end{equation}

\paragraph{Random witness maps.} The set of witness maps is the product
$\prod_{v:d_v\ge1}O_v$. A \emph{uniformly random witness map} is a uniform element of this
finite set; equivalently, the choices $\omega(v)$ are independent and $\omega(v)$ is uniform
on $O_v$. For $t\in X$, a variable $v$ and a group $h$, let
\[
\pi_{t,v,h}=\PP\bigl(h^\omega_t(v)=h\bigr)=
\begin{cases}
|\{e\ni v:\alpha_t(e)=h\}|/d_v&\text{if }d_v\ge1,\\
1\text{ if }v\in S_{t,h},\ 0\text{ otherwise}&\text{if }d_v=0,
\end{cases}
\]
and $\Nbar_{t,h}=\E N^\omega_{t,h}=\sum_v\pi_{t,v,h}$. Since $\sum_h\pi_{t,v,h}=1$,
$\sum_h\Nbar_{t,h}=n$. We drop the index $t$ when the point is fixed. The two constants of
the paper are
\[
c_0=\frac{1-e^{-2}}{4}=0.21616\ldots,\qquad
c_b=\frac{c_0}{2\sqrt2}\,\theta(b)=\frac{c_0\,(b-1)}{\sqrt2\,b^2}\quad(b\ge2).
\]

\begin{table}[t]
\centering
\small
\begin{tabular}{@{}>{\raggedright\arraybackslash}p{3.1cm}>{\raggedright\arraybackslash}p{9.6cm}l@{}}
\toprule
Symbol & Meaning & Source\\
\midrule
$n$, $M$, $b$ & number of variables, of occurrences; bound on occurrences per variable & \cite{DL}\\
$O_v$, $d_v$, $\nocc$ & occurrences of $v$; their number; number of $v$ with $d_v\ge1$ & here\\
$\Gr(F)$ & primal graph ($\mathrm{Pr}(F)$ in~\cite{DL}) & \cite{DL}\\
$P_o(T)$, $R$, $c_o(T)$ & deformation vector, degree bound, own coefficient & \cite{DL}\\
$X$, $q_t$, $\Lambda$ & evaluation points, Taylor orders, recovery functional & \cite{DL}\\
$S_{t,h}$, $\alpha_t$, $\cR_{t,h}$ & groups, clause assignment, measurement subspaces & \cite{DL}\\
$a_h$, $r_h$, $g_t$ & $|S_{t,h}|$, $\dim\cR_{t,h}$, number of groups at $t$ & here\\
$s_t$, $\rho_t$, $\mu_t$, $w_t$ & $\max_ha_h$, $\sum_hr_h$, number of measured groups, $\tw(\cI_t)$ & \cite{DL}\\
$\kappa(t)$ & cost charged at the point $t$ & \cite{DL}\\
$\omega$, $\Phi^\omega(T)$ & witness map, witness matrix & Def.~\ref{def:witness}\\
$h^\omega_t(v)$, $I^\omega_{t,h}$, $N^\omega_{t,h}$ & home group; variables with home group $h$; their number & Def.~\ref{def:def}\\
$\dfc^\omega_t$ & deficiency $\sum_h(a_h-N^\omega_{t,h})_+$ & Def.~\ref{def:def}\\
$\pi_{t,v,h}$, $\Nbar_{t,h}$ & $\PP(h^\omega_t(v)=h)$, $\E N^\omega_{t,h}$ for uniform $\omega$ & here\\
$\Sigma_t$, $h_t(v)$, $J_{t,h}$ & split variables; home group of a non-split variable; its fibres & Def.~\ref{def:split}\\
$\sep(F)$, $\theta(d)$, $\sepw(F)$ & separator number; weight $2(d-1)/d^2$; weighted separator number & Def.~\ref{def:sep}\\
$c_0$, $c_b$ & $(1-e^{-2})/4$; $c_0(b-1)/(\sqrt2b^2)$ & here\\
\bottomrule
\end{tabular}
\caption{Notation. The number of measured groups $\mu_t$ of~\cite{DL} appears only in
Section~\ref{sec:related} and Section~\ref{sec:dp}.}\label{tab:notation}
\end{table}

\section{Main results}\label{sec:results}

Throughout, $F$ is a CNF formula on $n\ge1$ variables and a \emph{scheme} is a clause-product
interpolation scheme for $F$ in the sense of Definition~5.1 of~\cite{DL}
(Section~\ref{sec:def}). No cost model is used before Corollary~\ref{cor:main}.

\begin{lemma}[Hermite counting]\label{lem:hermite}
In every scheme, $\sum_{t\in X}(q_t+1)\ge MR+1$. In particular $X\ne\emptyset$.
\end{lemma}

\begin{lemma}[Witness determinant]\label{lem:det}
For every scheme and every witness map $\omega$,
\[
\sum_{t\in X}\bigl(\dfc^\omega_t-\rho_t\bigr)_+\ \le\ \sum_{t\in X}\dim\ker\Phi^\omega(t)
\ \le\ \deg\det\Phi^\omega\ \le\ \nocc R\ \le\ \min(n,M)\,R .
\]
More precisely, for every $t\in X$,
$\dim\ker\Phi^\omega(t)\ge\sum_h\bigl(a_h-r_h-N^\omega_{t,h}\bigr)_+\ge\dfc^\omega_t-\rho_t$.
\end{lemma}

Figure~\ref{fig:blocks} shows the structure behind Lemma~\ref{lem:det}.

\begin{figure}[t]
\[
\Phi^\omega(t)\;=\;
\begin{array}{r|ccc}
 & \text{columns } S_{1} & \text{columns } S_{2} & \text{columns } S_{3}\\ \hline
\text{rows } I^\omega_{t,1} & \text{free} & \in\cR_{t,2} & \in\cR_{t,3}\\
\text{rows } I^\omega_{t,2} & \in\cR_{t,1} & \text{free} & \in\cR_{t,3}\\
\text{rows } I^\omega_{t,3} & \in\cR_{t,1} & \in\cR_{t,2} & \text{free}
\end{array}
\]
\caption{Block structure of the witness matrix at a point $t$, for three groups. Rows are
indexed by variables, sorted by home group; columns are sorted by group. By
condition~\eqref{eq:three}, a row whose home group is not $h$ has its restriction to the
columns $S_h$ inside $\cR_{t,h}$. The column block $S_h$ therefore has rank at most
$\min(a_h,\,N^\omega_{t,h}+r_h)$. Note that $N^\omega_{t,h}$, the number of rows with home
group $h$, need not equal $a_h=|S_h|$.}\label{fig:blocks}
\end{figure}

\begin{lemma}[Anticoncentration]\label{lem:anti}
Let $B$ be a sum of finitely many independent Bernoulli random variables, with variance $V$,
and let $a\in\mathbb R$. Then $\E|B-a|\ge c_0\min(V,\sqrt V)$, where $c_0=(1-e^{-2})/4$.
\end{lemma}

\begin{lemma}[Expected deficiency]\label{lem:expdef}
Let $\omega$ be a uniformly random witness map. For every scheme and every $t\in X$,
\[
\E\,\dfc^\omega_t\ \ge\ \max\Bigl(\frac12\sum_h\bigl|\Nbar_{t,h}-a_h\bigr|,\ \
\frac{c_0}{2\sqrt{2s_t}}\sum_{v\in\Sigma_t}\theta(d_v)\Bigr).
\]
\end{lemma}

\begin{lemma}[Split variables form a separator]\label{lem:split}
For every scheme and every $t\in X$, at least one of the following holds:
\begin{enumerate}
\item[\textup{(1)}] $|J_{t,h}|\le 2n/3$ for every group $h$. Then $\Sigma_t$ is a
 $2/3$-separator of $F$; in particular $|\Sigma_t|\ge\sep(F)$ and
 $\sum_{v\in\Sigma_t}\theta(d_v)\ge\sepw(F)$.
\item[\textup{(2)}] $|J_{t,h_0}|>2n/3$ for some group $h_0$. Then
 $\frac12\sum_h|\Nbar_{t,h}-a_h|>\frac{2n}{3}-s_t$.
\end{enumerate}
\end{lemma}

\begin{theorem}[Variable-uniform deformations: no absorption]\label{thm:uniform}
Suppose that the deformation is variable-uniform, that is, $P_{e,v}=P_{e',v}$ for all clauses
$e,e'$ containing $v$. Then some $t\in X$ satisfies
\[
\rho_t+q_t+1\ >\ \min\Bigl(\sep(F),\ \frac{2n}{3}-s_t\Bigr).
\]
\end{theorem}

\begin{theorem}[General deformations: at most a $\sqrt s$ factor]\label{thm:general}
Every scheme for $F$ has a point $t\in X$ with
\[
\rho_t+q_t+1\ >\ \min\Bigl(\frac{c_0}{2\sqrt2}\cdot\frac{\sepw(F)}{\sqrt{s_t}},\ \
\frac{2n}{3}-s_t\Bigr).
\]
This holds for arbitrary clause-dependent deformations, partitions, assignments $\alpha_t$
and subspaces $\cR_{t,h}$. If every variable occurs at most $b$ times, $b\ge2$, then some
$t\in X$ satisfies
\[
\rho_t+q_t+1\ >\ \min\Bigl(c_b\,\frac{\sep(F)}{\sqrt{s_t}},\ \frac{2n}{3}-s_t\Bigr),
\qquad c_b=\frac{c_0\,(b-1)}{\sqrt2\,b^2},\quad c_3=\frac{1-e^{-2}}{18\sqrt2}=0.03396\ldots
\]
\end{theorem}

The constant $c_0$ is not optimised (Remark~\ref{rem:const}).

\begin{lemma}[Random 3-CNF]\label{lem:random}
Fix $\Delta>3$. Let $F$ be a random 3-CNF on $n$ variables with $m=\Delta n$ independent
clauses ($\Delta n$ is assumed to be an integer), each on three distinct uniformly random variables with random signs. There are
constants $0<\eta'_\Delta\le\eta_\Delta\le1/9$ such that, with probability tending to $1$
as $n\to\infty$, $\sep(F)\ge\eta_\Delta n$ and $\sepw(F)\ge\eta'_\Delta n$.
\end{lemma}

This is the model of \cite[Appendix~A]{DL}; the proof adapts that of \cite[Lemma~8.2]{DL}.

\begin{corollary}\label{cor:main}
Let $F$ belong to one of the following families.
\begin{itemize}
\item[\textup{(i)}] Every variable occurs at most $b$ times and $\sep(F)\ge cn$, for fixed
 $b\ge2$ and $c>0$. Put $\gamma=c_bc$.
\item[\textup{(ii)}] $F$ is a random 3-CNF as in Lemma~\ref{lem:random}, on the event of
 that lemma (which has probability $1-o(1)$). Put
 $\gamma=\frac{c_0}{2\sqrt2}\eta'_\Delta$ and $c=\eta_\Delta$.
\end{itemize}
In both cases $\gamma\le1/6$ and $c\le1$. There is $n_0$, depending only on $\gamma$ and, in
part~1, on $\ell$, such that for $n\ge n_0$ the following hold for every scheme for $F$.
\begin{enumerate}
\item \textup{(The design of \cite[Section~9]{DL} does not exist.)} Let $\ell\ge0$ be fixed.
 If $s_t\le\sqrt n\log^{\ell}n$ for all $t\in X$, then some $t\in X$ has
 \[\rho_t+q_t+1>\gamma\,n^{3/4}/\log^{\ell/2}n .\]
 In particular, no scheme has $s_t,\ q_t,\ \rho_t\le\sqrt n\log^{\ell}n$ for all $t\in X$.
 This statement involves no cost model.
\item \textup{(Product-type cost.)} The product-type charge of the scheme is at least
 $2^{(\gamma n^{2/3}-1)/2}=2^{\Omega(n^{2/3})}$.
\item \textup{(Variable-uniform schemes.)} If the deformation is variable-uniform, the
 product-type charge is at least $2^{\min(n/6,\ (\min(c,1/2)\,n-1)/2)}=2^{\Omega(n)}$. In
 family~\textup{(i)}, this part does not need the bound $b$ on occurrences.
\end{enumerate}
\end{corollary}

In all three parts the scheme is applied to $F$ itself. For family~(ii) this means: to the
random formula, not to the bounded-occurrence formula $F^\star$ that the preprint derives
from it, and not to sparsification leaves.

\begin{remark}[Balanced designs]\label{rem:balanced}
Let $F$ be as in Corollary~\ref{cor:main} and let $t$ be the point given by
Theorem~\ref{thm:general}. Suppose that, at this particular $t$, $s_t\le n/2$, and that the
design has $g_t$ groups with $s_t\le Kn/g_t$ for some $K\ge1$. Then
$q_t+1\ge\frac{\gamma}{2\sqrt K}\sqrt{ng_t}$, or the average number of rows per group
satisfies $\rho_t/g_t>\frac{\gamma}{2K}\sqrt{s_t}$. The proof is in
Section~\ref{sec:proofcor}. The statement concerns one point $t$, not every point, and it
needs $s_t\le n/2$ so that the first term of the minimum in Theorem~\ref{thm:general} is the
smaller one.
\end{remark}

\begin{table}[t]
\centering
\small
\begin{tabular}{@{}>{\raggedright\arraybackslash}p{2.5cm}>{\raggedright\arraybackslash}p{3.7cm}>{\raggedright\arraybackslash}p{2.1cm}>{\raggedright\arraybackslash}p{6.2cm}@{}}
\toprule
Result & Scheme applied to & Model & Bound\\
\midrule
\cite[Thm~7.1(ii)]{DL} & any $F$ with $s_t\le 2n/3$ & none & $q_t+b\,\mu_ts_t\ge\sep(F)$ at some $t$ \newline(measured groups charged)\\[2pt]
\cite[Thm~8.1]{DL} & any $F$ & DP-type & cost $\ge2^{E}$, \newline $E=\min(\sqrt{(\tw+1)/2}-1,\,(\tw+1)/2,\,\lfloor n/2\rfloor)$\\[2pt]
\cite[Thm~8.3]{DL} & $F^\star$ and sparsification leaves, random 3-CNF, $\Delta>3$ & DP-type & cost $\ge2^{\Omega_\Delta(\sqrt n)}$ w.h.p.\\[2pt]
Prop.~\ref{prop:dp} & any $F$ & DP-type & a trivial scheme has charge $2^{\lceil\sqrt n\rceil}$ (upper bound)\\[2pt]
Thm~\ref{thm:uniform} & any $F$, variable-uniform & none & $\rho_t+q_t+1>\min(\sep(F),\,2n/3-s_t)$ \newline at some $t$\\[2pt]
Thm~\ref{thm:general} & any $F$ with $d_v\le b$ & none & $\rho_t+q_t+1>\min(c_b\sep(F)/\sqrt{s_t},\,2n/3-s_t)$ \newline at some $t$ (rows charged)\\[2pt]
Cor.~\ref{cor:main}(2) & $F$ itself, families (i), (ii) & product-type & cost $\ge2^{\Omega(n^{2/3})}$\\[2pt]
Cor.~\ref{cor:main}(3) & $F$ itself, variable-uniform & product-type & cost $\ge2^{\Omega(n)}$\\
\bottomrule
\end{tabular}
\caption{Comparison with the bounds of~\cite{DL}. The first line is our reading of
\cite[Thm~7.1(ii)]{DL} on separators (Section~\ref{sec:related}). ``DP-type'' bounds also hold
in the product-type model, which charges more. The lines for~\cite{DL} and for this paper
concern different formulas in the case of Theorem~8.3 ($F^\star$ and leaves there, $F$
here).}\label{tab:compare}
\end{table}

\section{Proofs}\label{sec:proofs}

\subsection{Proof of Lemma~\ref{lem:hermite} (Hermite counting)}

The argument is that of \cite[Lemma~6.2]{DL}, applied to a different polynomial. Suppose
$\sum_{t\in X}(q_t+1)\le MR$, and let $f(T)=\prod_{t\in X}(T-t)^{q_t+1}$ (so $f=1$ if
$X=\emptyset$). Then $\deg f\le MR$. For every $t\in X$ and $j\le q_t$ we have
$f^{(j)}(t)=0$, because $(T-t)^{q_t+1}$ divides $f$. So the Hermite data of $f$ are all zero
and, $\Lambda$ being linear, $\Lambda(f)=0$. But
$f(0)=\prod_{t\in X}(-t)^{q_t+1}\ne0$ because $0\notin X$. This contradicts the exactness
$\Lambda(f)=f(0)$ required in part~(b) of the definition. Hence
$\sum_{t\in X}(q_t+1)\ge MR+1\ge1$, and $X\ne\emptyset$. \qed

\subsection{Proof of Lemma~\ref{lem:det} (witness determinant)}

Fix a scheme and a witness map $\omega$, and write $\Phi=\Phi^\omega$.

\paragraph{Step 1: $\det\Phi$ is a nonzero polynomial of degree at most $\nocc R$.} If
$d_v\ge1$, row $v$ of $\Phi(0)$ is $P_{\omega(v)}(0)=e_v$, because $\omega(v)$ is an
occurrence of $v$. If $d_v=0$, row $v$ is $e_v$ by definition. So $\Phi(0)$ is the identity
matrix and $\det\Phi(0)=1$; in particular $\det\Phi\ne0$. By the Leibniz formula, $\det\Phi$
is a signed sum of products of $n$ entries, one in each row. The entries of a row $v$ with
$d_v\ge1$ have degree at most $R$, and the entries of a row with $d_v=0$ are constants. So
every product has degree at most $\nocc R$, and $\deg\det\Phi\le\nocc R\le\min(n,M)R$.

\paragraph{Step 2: rank bound at a point.} Fix $t\in X$ and a group $h$. Let $B_h$ be the
$n\times a_h$ matrix formed by the columns of $\Phi(t)$ indexed by $S_h$. Let $v$ be a
variable with $h^\omega_t(v)\ne h$. We claim that row $v$ of $B_h$ lies in $\cR_{t,h}$.
\begin{itemize}
\item If $d_v\ge1$ and $\omega(v)=(e,v)$, then $\alpha_t(e)=h^\omega_t(v)\ne h$, and row $v$
 of $B_h$ is $P_{e,v}(t)|_{S_h}$, which lies in $\cR_{t,h}$ by~\eqref{eq:three}.
\item If $d_v=0$, then $h^\omega_t(v)$ is the group containing $v$, so $v\notin S_h$, and
 row $v$ of $B_h$ is $e_v|_{S_h}=0\in\cR_{t,h}$.
\end{itemize}
So the row space of $B_h$ is contained in the sum of $\cR_{t,h}$ and the span of the
$N^\omega_{t,h}$ rows indexed by $I^\omega_{t,h}$. Since $B_h$ has $a_h$ columns,
\[
\rank B_h\ \le\ \min\bigl(a_h,\ r_h+N^\omega_{t,h}\bigr).
\]
The groups partition the columns of $\Phi(t)$, so the column space of $\Phi(t)$ is the sum
of the column spaces of the $B_h$, and $\rank\Phi(t)\le\sum_h\rank B_h$. With
$\sum_ha_h=n$ this gives
\[
\dim\ker\Phi(t)=n-\rank\Phi(t)\ \ge\ \sum_h\Bigl(a_h-\min\bigl(a_h,r_h+N^\omega_{t,h}\bigr)\Bigr)
=\sum_h\bigl(a_h-r_h-N^\omega_{t,h}\bigr)_+ .
\]
For real $y$ and $r\ge0$ we have $(y-r)_+\ge y_+-r$. Applying this with
$y=a_h-N^\omega_{t,h}$ and $r=r_h$ and summing over $h$,
\[
\dim\ker\Phi(t)\ \ge\ \sum_h\bigl(a_h-N^\omega_{t,h}\bigr)_+-\sum_hr_h=\dfc^\omega_t-\rho_t .
\]
As $\dim\ker\Phi(t)\ge0$, also $\dim\ker\Phi(t)\ge(\dfc^\omega_t-\rho_t)_+$.

\paragraph{Step 3: order of vanishing.} Fix $t\in X$ and let $k=\dim\ker\Phi(t)$. Choose an
invertible matrix $U\in\Q^{n\times n}$ whose first $k$ columns form a basis of
$\ker\Phi(t)$. The matrix $\Phi(T)U$ has entries in $\Q[T]$, and its first $k$ columns
vanish at $T=t$. So every entry of these columns is divisible by $T-t$. Factoring $T-t$ out
of each of these $k$ columns gives $\det(\Phi(T)U)=(T-t)^k\det M(T)$ for a matrix $M(T)$
over $\Q[T]$. Since $\det(\Phi U)=\det U\cdot\det\Phi$ and $\det U$ is a nonzero constant,
$(T-t)^{k}$ divides $\det\Phi$.

\paragraph{Step 4: conclusion.} The polynomials $(T-t)^{\dim\ker\Phi(t)}$, $t\in X$, are
pairwise coprime and each divides $\det\Phi$. So their product divides $\det\Phi$, which is
nonzero by Step~1. Comparing degrees,
$\sum_{t\in X}\dim\ker\Phi(t)\le\deg\det\Phi\le\nocc R$. Together with Step~2 this proves the
lemma. \qed

\begin{remark}\label{rem:mn}
Only $\deg\det\Phi^\omega\le MR$ is used below. This holds for every formula because
$\nocc\le M$; no assumption such as ``every variable occurs'' or $M\ge n$ is needed.
\end{remark}

\subsection{Proof of Lemma~\ref{lem:anti} (anticoncentration)}

Write $B=\sum_{v}\beta_v$ with independent $\beta_v\sim\mathrm{Bernoulli}(p_v)$, so
$V=\sum_vp_v(1-p_v)$. If $V=0$ there is nothing to prove. Let $B'=\sum_v\beta'_v$ be an
independent copy of $B$. By the triangle inequality $|B-B'|\le|B-a|+|B'-a|$, so
\begin{equation}\label{eq:sym}
\E|B-a|\ \ge\ \tfrac12\,\E|D|,\qquad D=B-B'=\sum_v\xi_v,\quad \xi_v=\beta_v-\beta'_v .
\end{equation}
The $\xi_v$ are independent, take values in $\{-1,0,1\}$, and are symmetric:
$\PP(\xi_v=1)=\PP(\xi_v=-1)=p_v(1-p_v)$. Put $\tau_v=\PP(\xi_v\ne0)=2p_v(1-p_v)$. Then
$\E\xi_v=0$, $\E\xi_v^2=\E\xi_v^4=\tau_v$, and $\sum_v\tau_v=2V$. In particular
$\E D^2=\sum_v\tau_v=2V$.

\paragraph{(i) A bound for all $V$.} Let $K=\{v:\xi_v\ne0\}$. Conditionally on $K$, the
values $(\xi_v)_{v\in K}$ are independent uniform signs, because the $\xi_v$ are independent
and symmetric. If $|K|=j\ge1$, then $\PP(D=0\mid K)$ is $0$ for odd $j$ and
$\binom{j}{j/2}2^{-j}\le\frac12$ for even $j$. Hence $\PP(D\ne0\mid K)\ge\frac12$ whenever
$K\ne\emptyset$. Also
\[
\PP(K\ne\emptyset)=1-\prod_v(1-\tau_v)\ \ge\ 1-e^{-\sum_v\tau_v}=1-e^{-2V}.
\]
Since $D$ is an integer, $|D|\ge1$ when $D\ne0$, so
$\E|D|\ge\PP(D\ne0)\ge\frac12(1-e^{-2V})$, and by~\eqref{eq:sym},
\begin{equation}\label{eq:smallV}
\E|B-a|\ \ge\ \tfrac14\bigl(1-e^{-2V}\bigr).
\end{equation}

\paragraph{(ii) A bound for $V\ge1$.} Expanding $D^4$ and using independence and
$\E\xi_v=\E\xi_v^3=0$, only the terms $\xi_v^4$ and $\xi_u^2\xi_v^2$ ($u\ne v$) have nonzero
expectation, and each unordered pair $\{u,v\}$ appears $\binom42=6$ times:
\[
\E D^4=\sum_v\tau_v+3\sum_{u\ne v}\tau_u\tau_v\ \le\ 2V+3(2V)^2=2V+12V^2 .
\]
By H\"older's inequality with exponents $3/2$ and $3$,
$\E D^2=\E\bigl(|D|^{2/3}|D|^{4/3}\bigr)\le(\E|D|)^{2/3}(\E D^4)^{1/3}$, hence
\[
\E|D|\ \ge\ \frac{(\E D^2)^{3/2}}{(\E D^4)^{1/2}}\ \ge\ \frac{(2V)^{3/2}}{(2V+12V^2)^{1/2}} .
\]
For $V\ge1$ we have $2V+12V^2\le14V^2$, so $\E|D|\ge\frac{2\sqrt2}{\sqrt{14}}\sqrt V
=\sqrt{4/7}\,\sqrt V$, and by~\eqref{eq:sym}
$\E|B-a|\ge\frac{1}{\sqrt7}\sqrt V\ge0.377\sqrt V\ge c_0\sqrt V$.

\paragraph{Conclusion.} The function $V\mapsto1-e^{-2V}$ is concave and vanishes at $0$, so
$(1-e^{-2V})/V$ is decreasing on $V>0$. For $0<V\le1$ this gives
$1-e^{-2V}\ge(1-e^{-2})V$, and~\eqref{eq:smallV} yields
$\E|B-a|\ge\frac{1-e^{-2}}{4}V=c_0V$. For $V\ge1$, (ii) gives $\E|B-a|\ge c_0\sqrt V$. Since
$\min(V,\sqrt V)$ equals $V$ for $V\le1$ and $\sqrt V$ for $V\ge1$, the lemma follows. \qed

\begin{remark}[On the constant]\label{rem:const}
The value of $c_0$ is forced by step~(i) at $V=1$; step~(ii) gives $0.377$. We made no
attempt to optimise it. In the numerical experiment of Section~\ref{sec:checks} the smallest
observed ratio $\E|B-a|/\min(V,\sqrt V)$ was $0.7397$.
\end{remark}

\subsection{Proof of Lemma~\ref{lem:expdef} (expected deficiency)}

Fix $t\in X$ and drop it from the notation; write $s=s_t$ and $N_h=N^\omega_{t,h}$. Recall
$\sum_hN_h=n=\sum_ha_h$. So $\sum_h(a_h-N_h)=0$, the positive and negative parts of the
numbers $a_h-N_h$ have the same sum, and
\[
\dfc^\omega_t=\sum_h(a_h-N_h)_+=\frac12\sum_h|N_h-a_h|,\qquad\text{hence}\qquad
\E\,\dfc^\omega_t=\frac12\sum_h\E|N_h-a_h| .
\]

\paragraph{First bound.} By Jensen's inequality,
$\E|N_h-a_h|\ge|\E N_h-a_h|=|\Nbar_h-a_h|$. Summing over $h$ gives the first term of the
maximum.

\paragraph{Structure of $N_h$.} For a fixed group $h$,
$N_h=\sum_v\mathbf 1[h^\omega_t(v)=h]$. The home group of $v$ depends only on $\omega(v)$
(and is deterministic if $d_v=0$), and the $\omega(v)$ are independent. So $N_h$ is a sum of
independent Bernoulli variables with parameters $\pi_{v,h}$, with mean $\Nbar_h$ and variance
\[
V_h=\sum_v\pi_{v,h}(1-\pi_{v,h})\ \le\ \Nbar_h .
\]

\paragraph{Second bound, one group.} We claim that
\begin{equation}\label{eq:onegroup}
\E|N_h-a_h|\ \ge\ \frac{c_0}{\sqrt{2s}}\,V_h\qquad\text{for every group }h.
\end{equation}
Note that $s\ge1$ and $a_h\le s$.

\emph{Case $\Nbar_h\le2s$.} Then $V_h\le2s$. If $V_h\le1$, then
$\min(V_h,\sqrt{V_h})=V_h\ge V_h/\sqrt{2s}$ because $\sqrt{2s}\ge1$. If $V_h\ge1$, then
$\min(V_h,\sqrt{V_h})=\sqrt{V_h}=V_h/\sqrt{V_h}\ge V_h/\sqrt{2s}$. In both subcases
Lemma~\ref{lem:anti}, applied to $B=N_h$ and $a=a_h$, gives~\eqref{eq:onegroup}.

\emph{Case $\Nbar_h>2s$.} Then $a_h\le s<\Nbar_h/2$, so
$\E|N_h-a_h|\ge\Nbar_h-a_h>\Nbar_h/2\ge V_h/2$. Since $s\ge1$,
$c_0/\sqrt{2s}\le c_0/\sqrt2<1/2$, and~\eqref{eq:onegroup} follows.

\paragraph{Summing over groups.} Since $\sum_h\pi_{v,h}=1$,
\[
\sum_hV_h=\sum_v\Bigl(1-\sum_h\pi_{v,h}^2\Bigr),
\]
and every term of the outer sum is nonnegative. Let $v\in\Sigma_t$ and $d=d_v\ge2$. Then
$\pi_{v,h}=k_h/d$, where $k_h$ is the number of clauses containing $v$ that are assigned to
$h$; the $k_h$ are nonnegative integers with sum $d$, and at least two of them are positive.
Let $k_{h_1}$ be a largest one; then $1\le k_{h_1}\le d-1$ and
$\sum_{h\ne h_1}k_h^2\le\bigl(\sum_{h\ne h_1}k_h\bigr)^2=(d-k_{h_1})^2$. The function
$k\mapsto k^2+(d-k)^2$ is convex, so on $1\le k\le d-1$ it is at most its value
$(d-1)^2+1$ at the endpoints. Therefore $\sum_hk_h^2\le(d-1)^2+1$ and
\[
1-\sum_h\pi_{v,h}^2\ \ge\ 1-\frac{(d-1)^2+1}{d^2}=\frac{2(d-1)}{d^2}=\theta(d_v).
\]
Keeping only the split variables, $\sum_hV_h\ge\sum_{v\in\Sigma_t}\theta(d_v)$.
With~\eqref{eq:onegroup},
\[
\E\,\dfc^\omega_t=\frac12\sum_h\E|N_h-a_h|\ \ge\ \frac{c_0}{2\sqrt{2s}}\sum_hV_h
\ \ge\ \frac{c_0}{2\sqrt{2s}}\sum_{v\in\Sigma_t}\theta(d_v),
\]
which is the second term of the maximum. \qed

\subsection{Proof of Lemma~\ref{lem:split} (split variables form a separator)}

Fix $t\in X$.

\paragraph{Components of $\Gr(F)-\Sigma_t$.} Let $u,v\notin\Sigma_t$ be adjacent in
$\Gr(F)$. They share a clause $e$, so $d_u,d_v\ge1$, and since neither is split,
$h_t(u)=\alpha_t(e)=h_t(v)$ by Definition~\ref{def:split}. Hence $h_t$ is constant on every
connected component of $\Gr(F)-\Sigma_t$: every component is contained in one set $J_{t,h}$.
(Variables with $d_v=0$ are isolated vertices of $\Gr(F)$; their home group is the group
that contains them, by the convention of Definition~\ref{def:split}.)

\paragraph{Case (1).} If $|J_{t,h}|\le2n/3$ for every $h$, every component of
$\Gr(F)-\Sigma_t$ has at most $2n/3$ vertices, so $\Sigma_t$ is a $2/3$-separator and
$|\Sigma_t|\ge\sep(F)$. Every split variable has $d_v\ge2$, so $\Sigma_t$ is admissible in
the definition of $\sepw(F)$ and $\sum_{v\in\Sigma_t}\theta(d_v)\ge\sepw(F)$.

\paragraph{Case (2).} Suppose $|J_{t,h_0}|>2n/3$. Let $v\in J_{t,h_0}$. If $d_v\ge1$, all
clauses containing $v$ are assigned to $h_0$, so $\pi_{v,h_0}=1$. If $d_v=0$, then
$v\in S_{t,h_0}$ and $\pi_{v,h_0}=1$ by definition. Hence
$\Nbar_{h_0}\ge|J_{t,h_0}|>2n/3$, while $a_{h_0}\le s_t$. Since
$\sum_h\Nbar_h=n=\sum_ha_h$, the numbers $\Nbar_h-a_h$ sum to zero, and
\[
\frac12\sum_h\bigl|\Nbar_h-a_h\bigr|=\sum_h\bigl(\Nbar_h-a_h\bigr)_+\ \ge\ \Nbar_{h_0}-a_{h_0}
\ >\ \frac{2n}{3}-s_t .
\]
The two cases are exhaustive. \qed

\subsection{Proof of Theorem~\ref{thm:uniform} (variable-uniform deformations)}

Assume $P_{e,v}=P_v$ for every clause $e\ni v$. Then $\Phi^\omega$ does not depend on the
witness map: its row $v$ is $P_v(T)$ if $d_v\ge1$ and $e_v$ if $d_v=0$. Write $\Phi$ for
this matrix. (The deformation then substitutes the linear form $\langle\Phi_v(T),x\rangle$
for $x_v$ in every clause.)

\paragraph{Rank bound.} Fix $t\in X$ and a group $h$, and let $B_h$ be the block of columns
$S_h$ of $\Phi(t)$. We claim that row $v$ of $B_h$ lies in $\cR_{t,h}$ for every
$v\notin J_{t,h}$. There are three cases.
\begin{itemize}
\item[(a)] $v\in\Sigma_t$. Then $v$ occurs in clauses assigned to two different groups, so
 it occurs in some clause $e$ with $\alpha_t(e)\ne h$.
\item[(b)] $v\notin\Sigma_t$, $d_v\ge1$ and $h_t(v)\ne h$. Then every clause $e\ni v$ has
 $\alpha_t(e)=h_t(v)\ne h$.
\item[(c)] $d_v=0$ and $h_t(v)\ne h$. Then $v\notin S_h$, and row $v$ of $B_h$ is
 $e_v|_{S_h}=0\in\cR_{t,h}$.
\end{itemize}
In cases (a) and (b), condition~\eqref{eq:three} applied to the occurrence $(e,v)$ gives
\[
P_v(t)|_{S_h}=P_{e,v}(t)|_{S_h}\in\cR_{t,h};
\]
this is where variable-uniformity is used. We now argue as in Step~2 of the proof of
Lemma~\ref{lem:det}, with $J_{t,h}$ in place of $I^\omega_{t,h}$. We get
$\rank B_h\le\min(a_h,r_h+|J_{t,h}|)$ and
\begin{equation}\label{eq:Dt}
\dim\ker\Phi(t)\ \ge\ \sum_h\bigl(a_h-r_h-|J_{t,h}|\bigr)_+\ \ge\ D_t-\rho_t,
\qquad D_t:=\sum_h\bigl(a_h-|J_{t,h}|\bigr)_+ .
\end{equation}

\paragraph{A good point.} Steps~1, 3 and~4 of the proof of Lemma~\ref{lem:det} apply to
$\Phi$. With~\eqref{eq:Dt} and Lemma~\ref{lem:hermite} they give
\[
\sum_{t\in X}(D_t-\rho_t)_+\ \le\ \sum_{t\in X}\dim\ker\Phi(t)\ \le\ \nocc R\ \le\ MR
\ <\ \sum_{t\in X}(q_t+1).
\]
So there is $t\in X$ with $(D_t-\rho_t)_+<q_t+1$, that is, $D_t<\rho_t+q_t+1$. Fix this $t$.

\paragraph{Lower bound on $D_t$.} Apply the case distinction of Lemma~\ref{lem:split}.

\emph{Case (1).} $\Sigma_t$ is a $2/3$-separator. Since the $J_{t,h}$ partition
$[n]\setminus\Sigma_t$,
\[
D_t\ \ge\ \sum_h\bigl(a_h-|J_{t,h}|\bigr)=n-(n-|\Sigma_t|)=|\Sigma_t|\ \ge\ \sep(F).
\]

\emph{Case (2).} $|J_{t,h_0}|>2n/3$ for some $h_0$. Then $\sum_{h\ne h_0}|J_{t,h}|<n/3$, and
\[
D_t=n-\sum_h\min\bigl(a_h,|J_{t,h}|\bigr)\ \ge\ n-a_{h_0}-\sum_{h\ne h_0}|J_{t,h}|
\ >\ \frac{2n}{3}-s_t .
\]
In both cases $\rho_t+q_t+1>D_t\ge\min(\sep(F),2n/3-s_t)$, with strict inequality. \qed

\subsection{Proof of Theorem~\ref{thm:general} (general deformations)}

\paragraph{Averaging the determinant lemma.} Let $\omega$ be a uniformly random witness map
and write $\E\,\dfc_t=\E\,\dfc^\omega_t$. For each $t$ the function $y\mapsto(y-\rho_t)_+$ is
convex, so by Jensen's inequality $(\E\,\dfc_t-\rho_t)_+\le\E(\dfc^\omega_t-\rho_t)_+$.
Summing over $t$, and using Lemma~\ref{lem:det} for every $\omega$, Remark~\ref{rem:mn} and
Lemma~\ref{lem:hermite},
\[
\sum_{t\in X}\bigl(\E\,\dfc_t-\rho_t\bigr)_+\ \le\ \E\sum_{t\in X}\bigl(\dfc^\omega_t-\rho_t\bigr)_+
\ \le\ \nocc R\ \le\ MR\ <\ \sum_{t\in X}(q_t+1).
\]
Hence there is a point $t\in X$ with $(\E\,\dfc_t-\rho_t)_+<q_t+1$, that is,
\begin{equation}\label{eq:goodpoint}
\E\,\dfc_t\ <\ \rho_t+q_t+1 .
\end{equation}

\paragraph{Lower bound at that point.} Apply Lemma~\ref{lem:split} at this $t$.

\emph{Case (1).} Then $\sum_{v\in\Sigma_t}\theta(d_v)\ge\sepw(F)$; in particular
$\sepw(F)<+\infty$. The second bound of Lemma~\ref{lem:expdef} gives
\[
\E\,\dfc_t\ \ge\ \frac{c_0}{2\sqrt{2s_t}}\,\sepw(F)=\frac{c_0}{2\sqrt2}\cdot\frac{\sepw(F)}{\sqrt{s_t}} .
\]

\emph{Case (2).} The first bound of Lemma~\ref{lem:expdef} and Lemma~\ref{lem:split}(2) give
$\E\,\dfc_t>2n/3-s_t$.

In both cases $\E\,\dfc_t$ is at least the minimum in the theorem; recall the convention
that this minimum is $2n/3-s_t$ when $\sepw(F)=+\infty$. With~\eqref{eq:goodpoint} this
gives the first statement.

\paragraph{Bounded occurrences.} Suppose $d_v\le b$ for all $v$, with $b\ge2$.
Then~\eqref{eq:sepw} gives $\sepw(F)\ge\theta(b)\sep(F)$, and
\[
\frac{c_0}{2\sqrt2}\cdot\frac{2(b-1)}{b^2}=c_b,\qquad
c_3=\frac{1-e^{-2}}{4}\cdot\frac{2}{9\sqrt2}=\frac{1-e^{-2}}{18\sqrt2}=0.03396\ldots
\]
This gives the second statement. \qed

\subsection{Proof of Lemma~\ref{lem:random} (random 3-CNF)}

We use two properties of the random formula, which are Lemmas~A.1 and~A.2 of~\cite{DL};
their proofs are reproduced in Appendix~\ref{app:random}. With probability $1-o(1)$:
\begin{itemize}
\item[(A1)] for every set $S$ of variables with $n/3\le|S|\le2n/3$, at least $m/3$ clauses
 contain both a variable of $S$ and a variable outside $S$;
\item[(A2)] for every $j\ge1$, every set of $j$ variables meets at most $C_\Delta j\log(en/j)$
 clauses, where $C_\Delta=3e^2\Delta$.
\end{itemize}
Assume that (A1) and (A2) hold.

\paragraph{Bound on $\sep(F)$.} Let $Y$ be a $2/3$-separator with $|Y|\le n/9$. The
components of $\Gr(F)-Y$ have at most $2n/3$ vertices each and at least $8n/9$ vertices in
total. Choose a union $A$ of components as follows. If some component has at least $n/3$
vertices, let $A$ be that component. Otherwise every component has fewer than $n/3$
vertices; add components one at a time until the union has at least $n/3$ vertices, and let
$A$ be this union, which has fewer than $2n/3$ vertices. In both cases $n/3\le|A|\le2n/3$.

Let $e$ be a clause that contains a variable of $A$ and a variable outside $A$. If $e$
contained no variable of $Y$, its variables would be pairwise adjacent in $\Gr(F)-Y$, hence
in one component, which is impossible because $A$ is a union of components. So every such
clause contains a variable of $Y$. By (A1) there are at least $m/3>0$ such clauses; in
particular $Y\ne\emptyset$. By (A2) with $j=|Y|\ge1$,
\[
\frac{\Delta n}{3}=\frac m3\ \le\ C_\Delta\,|Y|\log\frac{en}{|Y|} .
\]
With $y=|Y|/n\in(0,1/9]$ this reads $\Delta/(3C_\Delta)\le y\log(e/y)$. The function
$y\mapsto y\log(e/y)$ is increasing on $(0,1]$ and tends to $0$ as $y\to0$. So there is a
constant $y^*_\Delta\in(0,1/9]$ such that $y\log(e/y)<\Delta/(3C_\Delta)$ for all
$0<y<y^*_\Delta$, and therefore $y\ge y^*_\Delta$. Put $\eta_\Delta=y^*_\Delta\le1/9$. We
have shown that every $2/3$-separator has more than $n/9$ or at least $\eta_\Delta n$
vertices. Hence $\sep(F)\ge\eta_\Delta n$.

\paragraph{Bound on $\sepw(F)$.} Every clause has exactly three distinct variables, so
$\sum_vd_v=M=3\Delta n$, and for every integer $D\ge1$ at most $3\Delta n/D$ variables have
$d_v>D$. Take $D=\lceil6\Delta/\eta_\Delta\rceil\ge2$, so that $3\Delta n/D\le\eta_\Delta n/2$.
Let $Y$ be a $2/3$-separator with $d_v\ge2$ for all $v\in Y$. Then $|Y|\ge\eta_\Delta n$, so
$Y$ contains at least $\eta_\Delta n/2$ variables with $2\le d_v\le D$, and each of them has
$\theta(d_v)\ge\theta(D)$ because $\theta$ is decreasing on $d\ge2$. Hence
$\sum_{v\in Y}\theta(d_v)\ge\frac{\eta_\Delta}{2}\theta(D)\,n$. Put
$\eta'_\Delta=\frac{\eta_\Delta}{2}\theta(D)$. Since $\theta\le\frac12$ on $d\ge2$,
$\eta'_\Delta\le\eta_\Delta/4$. If there is no such $Y$, $\sepw(F)=+\infty$. In both cases
$\sepw(F)\ge\eta'_\Delta n$. \qed

\subsection{Proof of Corollary~\ref{cor:main} and of Remark~\ref{rem:balanced}}
\label{sec:proofcor}

\paragraph{A common inequality.} In family~(i), the second statement of
Theorem~\ref{thm:general} and $\sep(F)\ge cn$ give a point $t\in X$ with
\begin{equation}\label{eq:star}
\rho_t+q_t+1\ >\ \min\Bigl(\frac{\gamma n}{\sqrt{s_t}},\ \frac{2n}{3}-s_t\Bigr),
\end{equation}
with $\gamma=c_bc$. In family~(ii), the first statement of Theorem~\ref{thm:general} and
$\sepw(F)\ge\eta'_\Delta n$ give~\eqref{eq:star} with
$\gamma=\frac{c_0}{2\sqrt2}\eta'_\Delta$. In family~(i), $c\le1$ because $\sep(F)\le n$, and
$c_b\le c_0/(4\sqrt2)<0.04$ because $(b-1)/b^2\le1/4$. In family~(ii),
$\eta'_\Delta\le\eta_\Delta\le1/9$. So $\gamma\le1/6$ and $c\le1$ in both cases.

\paragraph{Part 1.} Suppose $s_t\le\sqrt n\log^\ell n$ for all $t\in X$, and let $t$ be the
point of~\eqref{eq:star}. Then
$\gamma n/\sqrt{s_t}\ge\gamma n^{3/4}/\log^{\ell/2}n$ and
$2n/3-s_t\ge2n/3-\sqrt n\log^\ell n$. For $n$ large (depending on $\gamma$ and $\ell$) the
second quantity is at least $n/2\ge\gamma n^{3/4}/\log^{\ell/2}n$, so
$\rho_t+q_t+1>\gamma n^{3/4}/\log^{\ell/2}n$. For $n$ large this exceeds
$2\sqrt n\log^\ell n+1$, so $\rho_t$ and $q_t$ cannot both be at most $\sqrt n\log^\ell n$.

\paragraph{Part 2.} For $1\le j\le n/2$ we have
$\binom nj=\prod_{i=0}^{j-1}\frac{n-i}{j-i}\ge(n/j)^j\ge2^j$, and $\binom n0=2^0$. So, at
every point $t$, the product-type charge is at least
\begin{equation}\label{eq:charge}
\max\bigl(2^{s_t},\ 2^{\rho_t},\ 2^{\min(q_t,\lfloor n/2\rfloor)}\bigr).
\end{equation}
Let $t$ be the point of~\eqref{eq:star} and let $n\ge216$. If $s_t\ge n^{2/3}$, the charge
is at least $2^{n^{2/3}}$. Otherwise $s_t<n^{2/3}\le n/6$, so $2n/3-s_t>n/2$ and
$\gamma n/\sqrt{s_t}>\gamma n^{2/3}$; as $\gamma n^{2/3}\le n/2$, \eqref{eq:star} gives
$\rho_t+q_t+1>\gamma n^{2/3}$, hence $\max(\rho_t,q_t)>(\gamma n^{2/3}-1)/2$. If
$\rho_t>(\gamma n^{2/3}-1)/2$, the charge is at least $2^{\rho_t}$. If
$q_t>(\gamma n^{2/3}-1)/2$, then, since $(\gamma n^{2/3}-1)/2\le\lfloor n/2\rfloor$, the
charge is at least $2^{\min(q_t,\lfloor n/2\rfloor)}\ge2^{(\gamma n^{2/3}-1)/2}$. As
$\gamma<1$, in all cases the charge is at least $2^{(\gamma n^{2/3}-1)/2}$.

\paragraph{Part 3.} Theorem~\ref{thm:uniform} gives a point $t$ with
$\rho_t+q_t+1>\min(\sep(F),2n/3-s_t)$, and $\sep(F)\ge cn$ in both families (by
Lemma~\ref{lem:random} in family~(ii)); the bound $b$ is not used. If $s_t\ge n/6$, the
charge~\eqref{eq:charge} is at least $2^{n/6}$. Otherwise $2n/3-s_t>n/2$, so
$\rho_t+q_t+1>c'n$ with $c'=\min(c,1/2)$, and $\max(\rho_t,q_t)>(c'n-1)/2$. Since
$(c'n-1)/2\le\lfloor n/2\rfloor$, \eqref{eq:charge} is at least $2^{(c'n-1)/2}$, as in
Part~2. So the charge is at least $2^{\min(n/6,(c'n-1)/2)}$.

\paragraph{Remark~\ref{rem:balanced}.} Let $t$ be the point of~\eqref{eq:star} and suppose
$s_t\le n/2$ and $s_t\le Kn/g_t$. Since $\gamma\le1/6$ and $s_t\ge1$,
$\gamma n/\sqrt{s_t}\le n/6\le2n/3-s_t$, so~\eqref{eq:star} reads
$\rho_t+q_t+1>\gamma n/\sqrt{s_t}\ge\frac{\gamma}{\sqrt K}\sqrt{ng_t}$. If
$q_t+1<\frac{\gamma}{2\sqrt K}\sqrt{ng_t}$, then
$\rho_t>\frac{\gamma}{2\sqrt K}\sqrt{ng_t}$, and dividing by $g_t$,
$\rho_t/g_t>\frac{\gamma}{2\sqrt K}\sqrt{n/g_t}\ge\frac{\gamma}{2\sqrt K}\sqrt{s_t/K}
=\frac{\gamma}{2K}\sqrt{s_t}$. \qed

\section{The DP-type model is degenerate}\label{sec:dp}

The DP-type model of \cite[Section~5.2]{DL} charges $2^{s_t}$, $2^{w_t}$ and the Taylor
factor, but not $2^{\rho_t}$ (Section~\ref{sec:cost}). The following observation shows that
its charge is small on every formula, for a trivial scheme.

\begin{proposition}[A trivial scheme in the DP-type model]\label{prop:dp}
Let $F$ be any CNF formula on $n\ge1$ variables and let $R\ge0$. There is a scheme for $F$ in
the sense of Definition~5.1 of~\cite{DL}, with degree bound $R$, such that at every point
$t\in X$
\[
s_t=\lceil\sqrt n\,\rceil,\qquad w_t\le\lceil\sqrt n\,\rceil-1,\qquad q_t=0,\qquad\rho_t=n .
\]
Its DP-type charge is exactly $2^{\lceil\sqrt n\,\rceil}$, and its product-type charge is
$2^n$.
\end{proposition}

\begin{proof}
\emph{(a) Deformation.} Put $P_o(T)=e_v$, a constant vector, for every occurrence $o=(e,v)$.
Then $\deg P_o=0\le R$ and $P_o(0)=e_v$. The polynomial $H$ of~\eqref{eq:H} is the constant
$\#F$, and $\deg H\le MR$.

\emph{(b) Recovery.} Let $X=\{1,2,\dots,MR+1\}\subset\Q\setminus\{0\}$ and $q_t=0$ for all
$t\in X$. The Hermite data of $f$ are the values $(f(t))_{t\in X}$. Let
$\Lambda\bigl((y_t)_{t\in X}\bigr)=\sum_{t\in X}y_t\,L_t(0)$, where
$L_t(T)=\prod_{u\in X\setminus\{t\}}\frac{T-u}{t-u}$ is the Lagrange basis. This is a linear
functional, and $\Lambda(f)=f(0)$ for every polynomial $f$ of degree at most $|X|-1=MR$, by
Lagrange interpolation.

\emph{(c) Designs.} Let $s=\lceil\sqrt n\,\rceil$ and $g=\lceil n/s\rceil$. At every $t\in X$
use the same design: the groups $S_1,\dots,S_g$ are consecutive blocks of $[n]$, the first
$g-1$ of size $s$ and the last of size at most $s$; $\alpha_t(e)$ is an arbitrary group; and
$\cR_{t,h}=\Q^{S_h}$ for every $h$. Condition~\eqref{eq:three} holds because every vector of
$\Q^{S_h}$ lies in $\cR_{t,h}$.

\emph{Parameters.} Since $1\le s\le n$, the first group has exactly $s$ elements, so $s_t=s$.
Since $n/s\le n/\sqrt n=\sqrt n\le s$ and $s$ is an integer, $g\le s$. The interaction graph
$\cI_t$ has $g$ vertices, so $w_t=\tw(\cI_t)\le g-1\le s-1$. Finally
$\rho_t=\sum_h|S_h|=n$ and $\binom{n}{\min(q_t,\lfloor n/2\rfloor)}=\binom n0=1$. The DP-type
charge is $\max(2^{s},2^{w_t},1)=2^{s}$, and the product-type charge is
$\max(2^s,2^{n})=2^n$.
\end{proof}

\paragraph{Consequences.}
\begin{enumerate}
\item \emph{The bound of~\cite{DL} is tight in the DP-type model, for a trivial reason.} For
 every formula, the smallest DP-type charge of a scheme is at most
 $2^{\lceil\sqrt n\,\rceil}$. With \cite[Theorem~8.1]{DL}, on any family with primal
 treewidth at least $cn$ it is $2^{\Theta(\sqrt n)}$: for $n$ large the exponent lies
 between $\sqrt{(cn+1)/2}-1$ and $\lceil\sqrt n\,\rceil$. The theorems of~\cite{DL} are lower
 bounds and remain correct. What Proposition~\ref{prop:dp} shows is that no argument, by any
 technique, can prove more than $2^{\lceil\sqrt n\,\rceil}$ in that model. The sentence of
 \cite[Section~9]{DL} that ``our lower-bound technique cannot exceed
 $2^{\tilde O(\sqrt n)}$'' is therefore a property of the model and not only of the
 technique.
\item \emph{The gap of \cite[Section~9]{DL} is not open in the DP-type model.} The cost
 target $2^{\tilde O(\sqrt n)}$ ``from above'' is met there by the trivial scheme, with
 $\sqrt n$ rows per group instead of polylogarithmically many. The question whether
 measurements can absorb the cut is meaningful only in a model that charges rows. The
 product-type model does, and there Corollary~\ref{cor:main} applies; the trivial scheme has
 product-type charge $2^n$, consistently with Theorem~\ref{thm:uniform}.
\item \emph{Why the model under-charges.} The charge $2^{w_t}$ counts two values per measured
 group, whatever its number of rows; this is the justification quoted in
 Section~\ref{sec:cost}. In the trivial scheme a group is fully measured, and its reading
 determines its whole assignment, so it takes $2^{|S_h|}$ values. A dynamic programme over
 readings would then be a dynamic programme over assignments of blocks of variables, which
 the charge $2^{w_t}$ does not price. A repaired DP-type model must weight a group by its
 number of rows; see Section~\ref{sec:concl}.
\end{enumerate}

\section{Machine verification}\label{sec:lean}

The deterministic part of the argument is formalised in Lean~4~\cite{Lean4} with
Mathlib~\cite{Mathlib}. The sources are in the directory \texttt{lean/} of the repository
\begin{center}
\url{https://github.com/luluddof/pnp-recherche}
\end{center}
The toolchain is Lean \texttt{v4.37.0-rc1}, and the Mathlib commit is \texttt{e62b985f}, as
pinned in the file \texttt{lake-manifest.json}. The four source files \texttt{WitnessDet},
\texttt{Scheme}, \texttt{TheoremB} and \texttt{TheoremA} (in \texttt{lean/PnpLean/Article1/}) contain
no \texttt{sorry} and declare no axiom. To reproduce:
\begin{verbatim}
  cd lean && lake exe cache get && lake build PnpLean.Article1.TheoremA
\end{verbatim}
The target \texttt{PnpLean.Article1.TheoremA} imports the other three files. We ran the build
command with these versions (the Mathlib cache being already present) and it completed
successfully. We then ran the command \texttt{\#print axioms} on ten declarations of
Table~\ref{tab:lean}:
\begin{center}
\texttt{sum\_nullity\_le}, \texttt{deficiency\_le\_nullity\_add}, \texttt{lemma1},
\texttt{factR}, \texttt{avg\_lemma1},\\ \texttt{exists\_point}, \texttt{lemma3},
\texttt{theoremB}, \texttt{theoremB\_bounded}, \texttt{theoremA}.
\end{center}
For each of them it reported exactly the three standard axioms \texttt{propext},
\texttt{Classical.choice} and \texttt{Quot.sound}.

\begin{table}[tb]
\centering
\small
\begin{tabular}{@{}>{\raggedright\arraybackslash}p{3.6cm}>{\raggedright\arraybackslash}p{6.1cm}>{\raggedright\arraybackslash}p{4.9cm}@{}}
\toprule
Statement of this paper & Lean declaration & Status\\
\midrule
Lemma~\ref{lem:hermite} (Hermite counting) & \texttt{Scheme.factR} & formalised\\[2pt]
Lemma~\ref{lem:det}, Step 1 & \texttt{det\_ne\_zero\_of\_eval\_zero}, \newline\texttt{natDegree\_det\_le} & formalised with the bound $nR$ (not $\nocc R$)\\[2pt]
Lemma~\ref{lem:det}, Step 2 & \texttt{deficiency\_le\_nullity\_add}, \newline\texttt{Scheme.deficiency\_sub\_le\_nullity} & formalised\\[2pt]
Lemma~\ref{lem:det}, Steps 3--4 & \texttt{pow\_dvd\_det\_of\_nullity}, \newline\texttt{sum\_nullity\_le} & formalised\\[2pt]
Lemma~\ref{lem:det} & \texttt{Scheme.lemma1} & formalised with the bound $nR$\\[2pt]
Averaging step of Theorem~\ref{thm:general}, and~\eqref{eq:goodpoint} & \texttt{Scheme.avg\_lemma1}, \newline\texttt{Scheme.exists\_point} & formalised; \texttt{exists\_point} assumes $n\le M$\\[2pt]
Lemma~\ref{lem:split} & \texttt{Scheme.lemma3} & formalised, as the disjunction ``$\Sigma_t$ is a $2/3$-separator, or $\frac12\sum_h|\Nbar_{t,h}-a_h|>2n/3-s_t$''\\[2pt]
Theorem~\ref{thm:uniform} & \texttt{Scheme.theoremA} & formalised, unconditionally, assuming $n\le M$\\[2pt]
Theorem~\ref{thm:general} & \texttt{Scheme.theoremB}, \newline\texttt{Scheme.theoremB\_bounded} & formalised \emph{conditionally}: Lemma~\ref{lem:expdef} is an explicit hypothesis \texttt{Lemma2Hyp}; assumes $n\le M$\\
\midrule
Lemma~\ref{lem:anti}, Lemma~\ref{lem:expdef} & --- & \textbf{not formalised} (paper proof only)\\[2pt]
Uniform average $=$ independent uniform choices & --- & \textbf{not formalised} (argued on paper, Section~\ref{sec:ournot})\\[2pt]
Bound $\nocc R\le MR$ without $n\le M$ (Remark~\ref{rem:mn}) & --- & \textbf{not formalised}\\[2pt]
Lemma~\ref{lem:random}, Corollary~\ref{cor:main}, Remark~\ref{rem:balanced}, Proposition~\ref{prop:dp}, Appendix~\ref{app:random} & --- & \textbf{not formalised}\\
\bottomrule
\end{tabular}
\caption{What is and is not machine-checked. Declarations without the prefix
\texttt{Scheme.} are in the namespace \texttt{WitnessDet}; the others are in
\texttt{ClauseScheme.Scheme}. The Lean names \texttt{lemma1}, \texttt{lemma3},
\texttt{theoremA}, \texttt{theoremB}, \texttt{factR} follow an earlier numbering of the
results.}\label{tab:lean}
\end{table}


\paragraph{What the Lean theorems say.} In \texttt{Scheme.theoremB}, the expectation
$\E\,\dfc_t$ is defined as the arithmetic mean of $\dfc^\omega_t$ over the finite type of all
witness maps, and \texttt{Lemma2Hyp c0} is the statement of Lemma~\ref{lem:expdef} for that
mean, with a constant $c_0\ge0$ left as a parameter. The theorem states: if $n\le M$, if
\texttt{Lemma2Hyp c0} holds, and if $W$ is a real number such that
$W\le\sum_{v\in Y}\theta(d_v)$ for every $2/3$-separator $Y$ whose variables all have
$d_v\ge2$, then some $t\in X$ has
$\min\bigl(\frac{c_0}{2\sqrt2}W/\sqrt{s_t},\,2n/3-s_t\bigr)<\rho_t+q_t+1$. So the machine
check covers the logical skeleton of Theorem~\ref{thm:general} (determinant, Hermite
counting, averaging, separator lemma, case analysis and constants) and \emph{not} its
probabilistic ingredient. \texttt{Scheme.theoremA} states Theorem~\ref{thm:uniform} with no
unproved hypothesis other than $n\le M$.

\paragraph{Differences between the Lean definitions and those of the paper.}
\begin{enumerate}
\item \emph{Formula.} A formula is given only by its clause variable sets
 $I:\kappa\to\mathrm{Finset}(\mathrm{Fin}\ n)$ for a finite type $\kappa$ of clauses. Signs,
 the polynomials $\varphi_e$ and $H$ do not appear. The structure \texttt{Scheme} records
 what the proofs use: the vectors $P_o$ with entries of degree at most $R$ and
 $P_o(0)=e_v$; a set $X$ of rationals not containing $0$; orders $q$; and a linear functional
 $\Lambda$ on the Hermite data, exact on every polynomial of degree at most $MR$. This is a
 weaker set of hypotheses than Definition~5.1, so the Lean theorems apply to every scheme of
 the paper.
\item \emph{Groups.} The groups are indexed by one fixed finite type $G$ at all points, and
 the partition at $t$ is a map from variables to $G$; empty groups are allowed. A scheme of
 the paper is encoded by taking $G$ large enough and giving unused labels the empty group
 and the zero subspace. The design data are functions of all $t\in\Q$, but
 condition~\eqref{eq:three} is required only for $t\in X$.
\item \emph{Measurement subspaces.} $\cR_{t,h}$ is a subspace of $\Q^n$ whose vectors vanish
 outside $S_{t,h}$, instead of a subspace of $\Q^{S_{t,h}}$; condition~\eqref{eq:three} is
 stated for the vector $P_{e,v}(t)$ with its coordinates outside $S_{t,h}$ set to zero.
\item \emph{Separators.} \texttt{IsSep I Y} says that for every $v\notin Y$, the set of
 vertices reachable from $v$ in the primal graph with all edges touching $Y$ removed has at
 most $2n/3$ elements. This is equivalent to Definition~\ref{def:sep}. The numbers
 $\sep(F)$ and $\sepw(F)$ are not defined; the theorems take an arbitrary lower bound
 (\texttt{sepLB}, respectively $W$) valid for every separator, which gives the statements
 of the paper on taking the best lower bound.
\item \emph{Hypothesis $n\le M$.} The Lean statements of \texttt{exists\_point},
 \texttt{theoremA}, \texttt{theoremB} and \texttt{theoremB\_bounded} assume $n\le M$ and use
 $\deg\det\Phi^\omega\le nR\le MR$. The paper removes this hypothesis through the sharper
 bound $\deg\det\Phi^\omega\le\nocc R$ (Step~1 of the proof of Lemma~\ref{lem:det}); this
 one-line sharpening is not formalised.
\item \emph{Arithmetic.} Deficiencies and positive parts are truncated subtractions of
 natural numbers; $s_t$ is the supremum of $a_h$ over $G$; in \texttt{theoremB\_bounded} the
 bound on occurrences satisfies $b\ge2$.
\end{enumerate}
The correspondence between the Lean statements and the statements of the paper was checked
by reading, by AI agents and not by a human expert; a mismatch between a Lean definition and
its intended meaning is the main residual risk of this section.

\section{Computational checks}\label{sec:checks}

We ran four small scripts. They test the algebra and the inequalities on small instances.
\textbf{Such checks can only detect errors; they prove nothing, and in particular nothing
about the asymptotic statements.} The scripts are in the repository cited in
Section~\ref{sec:lean}, in \texttt{code/gen6/g6a/} and \texttt{code/gen6/g6r/}. We report
only what we observed when re-running unmodified copies of them for this version, with
Python~3.13.18, sympy~1.14.0 and networkx~3.7, on a Linux machine (single thread).
Table~\ref{tab:checks} gives the commands, run from the repository root.

\begin{table}[htbp]
\centering
\small
\begin{tabular}{@{}llr@{}}
\toprule
Check & Command & Time\\
\midrule
C1 & \texttt{cd code/gen6/g6a \&\& python3 run\_checks.py theorem} & 6 s\\
C2 & \texttt{cd code/gen6/g6a \&\& python3 run\_checks.py e2e} & 19 s\\
C3 & \texttt{cd code/gen6/g6a \&\& python3 anticonc\_check.py} & 0.5 s\\
C4 & \texttt{cd code/gen6/g6r \&\& python3 lemma1\_check.py} & 125 s\\
\bottomrule
\end{tabular}
\caption{Commands for the computational checks. C1 needs networkx; C2 and C4 need
sympy. All scripts use fixed seeds.}\label{tab:checks}
\end{table}

\paragraph{C1: Lemma~\ref{lem:det} and the bound of Lemma~\ref{lem:expdef}.}
\emph{Instances.} Random formulas with $n$ variables and $n$ clauses in which every variable
occurs in exactly three clauses, for $(n,|X|,R)=(24,8,2)$, $(30,10,2)$ and $(36,9,3)$, with
$X=\{1,\dots,|X|\}$ and groups of size $6$ in a random partition at each point. Four families
of designs: cut-and-measure with a variable-uniform deformation; cut-and-measure with a
clause-dependent deformation; a ``global-sum'' design; a random dense deformation with the
smallest admissible $\cR_{t,h}$. This makes $12$ runs.

\emph{What is asserted.} Condition~\eqref{eq:three}, exactly over $\Q$. For four random
witness maps per run: $\dim\ker\Phi^\omega(t)\ge\dfc^\omega_t-\rho_t$ at every point, and
$\sum_t\dim\ker\Phi^\omega(t)\le nR$, with exact ranks over $\Q$. At every point, that
$\E\,\dfc_t$ is at least the bound of Lemma~\ref{lem:expdef} with $c_0$ replaced by $0.216$;
here $\E\,\dfc_t$ is computed from the exact distribution of each $N_h$ in floating-point
arithmetic. Finally $\sum_t(\E\,\dfc_t-\rho_t)_+\le nR$.

\emph{What we observed.} All assertions passed in the $12$ runs. In the order of the four families, the largest value of
$\sum_t\dim\ker\Phi^\omega(t)$ over the four witness maps was $48$, $48$, $40$, $0$ for
$n=24$ (where $nR=48$); $60$, $60$, $50$, $0$ for $n=30$ (where $nR=60$); and $108$, $108$,
$99$, $0$ for $n=36$ (where $nR=108$). So the bound $nR$ of Lemma~\ref{lem:det} was attained by the
cut-and-measure designs. Over all points of all runs, $\E\,\dfc_t$ ranged from $2.57$ to
$8.46$ and $\rho_t$ from $10$ to $36$.

\emph{Limits of this check.} In every one of the $12$ runs, $\rho_t$ exceeded $\E\,\dfc_t$
at every point (by at least $6.9$), so $\sum_t(\E\,\dfc_t-\rho_t)_+=0$. \textbf{The
averaging step of Theorem~\ref{thm:general} is therefore never really exercised}: the
inequality $\sum_t(\E\,\dfc_t-\rho_t)_+\le nR$ holds trivially in these runs. These designs
have $\rho_t$ between $0.4n$ and $n$ and are far from the regime of
Theorem~\ref{thm:general}; they test neither the tightness of the $\sqrt s$ factor nor the
final inequalities of Theorems~\ref{thm:uniform} and~\ref{thm:general}. A stored output
file in the repository (\texttt{theorem\_checks.txt}) contains only the configurations
$n=24$ and $n=30$; in our run the configuration $n=36$ completed normally, and the output
for $n=24$ and $n=30$ was identical to the stored file.

\paragraph{C2: end-to-end recovery.} One formula with $n=12$ (three occurrences per
variable), three groups of four variables, $R=1$, $X=\{1,\dots,6\}$, $q_t=6$ (so
$\sum_t(q_t+1)=42>MR=36$), and two design families. The script computes $H(T)$ by brute
force modulo the prime $2^{61}-1$, recovers $H(0)$ from the Hermite data, and at each point
compares a product-type evaluation (enumeration of tuples of measurement readings, counting
inside groups) with the brute-force value of $H(t)$ in exact rational arithmetic.
\emph{Observed:} $\#F=977$; $H(0)=977$ and the Hermite-recovered value is $977$ for both
designs; the product-type evaluation equals $H(t)$ at all $6$ points for both designs, with
between $64$ and $3072$ tuples of readings. This checks the code and the meaning of
Definition~5.1 on one instance; it does not test our theorems.

\paragraph{C3: anticoncentration.} $4000$ random sums of at most $40$ independent Bernoulli
variables; for each, the minimum over integers $a$ of $\E|B-a|$ divided by
$\min(V,\sqrt V)$, in floating-point arithmetic. \emph{Observed:} the smallest ratio is
$0.7397$, against $c_0\approx0.2162$ in Lemma~\ref{lem:anti}. This is a sanity check on
random instances, not a proof and not a search for worst cases.

\paragraph{C4: an independent test of Lemma~\ref{lem:det}.} The script builds random schemes
with $n=9$, $7$ clauses of three variables, $R=2$, $4$ points, $3$ groups and one random
measurement row per group, by solving the linear constraints~\eqref{eq:three} for the
coefficients of the deformation (sympy, exact arithmetic). For each scheme and six random
witness maps it asserts $P_o(0)=e_v$, $\det\Phi^\omega(0)=1$, $\deg\det\Phi^\omega\le nR$,
$\dim\ker\Phi^\omega(t)\ge\dfc^\omega_t-\rho_t$, that the multiplicity of the root $t$ of
$\det\Phi^\omega$ is at least $\dim\ker\Phi^\omega(t)$, and
$\sum_t(\dfc^\omega_t-\rho_t)_+\le nR$. \emph{Observed:} of the six seeds $0,\dots,5$, three
(seeds $0$, $4$, $5$) produced no scheme, because the linear system was infeasible, and the
script returned nothing for them. For seeds $1$, $2$, $3$ all assertions passed, on
$3\times6=18$ witness maps in total. In these three schemes $\rho_t=3$ at every point, and
the largest value of $\sum_t(\dfc^\omega_t-\rho_t)_+$ was $1$, $3$ and $1$, against $nR=18$.
The test is therefore weak: it is far from the equality case.

\paragraph{What is not checked numerically.} Lemma~\ref{lem:split},
Lemma~\ref{lem:random}, Corollary~\ref{cor:main} and Proposition~\ref{prop:dp} were not
checked numerically for this version, and no run exhibits a point with
$\E\,\dfc_t>\rho_t$.

\section{Methodology and use of AI tools}\label{sec:ai}

We state plainly how this work was produced.

\paragraph{AI-generated mathematics.} The results of this paper were found, proved, written
up and reviewed by AI agents, instances of Claude (Anthropic), working under the direction of
the author, who supervised the workflow and decided what to publish. The author did not
write the proofs.

\paragraph{Reviews, and their limits.} A first agent produced the determinant argument,
proofs and test scripts. A second agent carried out an adversarial review of each lemma and
wrote the independent test C4. A third agent wrote a referee report on an earlier version of
this text; that report found the proofs correct as far as it could check, and required the
changes of scope that are now in the paper: the DP-type model is degenerate rather than
``open'' (Section~\ref{sec:dp}), the comparison is with \cite[Theorem~8.1 and Lemma~8.2]{DL}
and not with their Theorem~8.3, and the computational section was corrected. This version
was then written by a further agent, which re-derived the proofs and re-ran the checks.
\textbf{These reviews are not independent of the provers}: all agents are instances of the
same family of models and may share blind spots. They do not replace review by human
experts.

\paragraph{Machine verification.} To reduce this risk, the deterministic part of the
argument was formalised in Lean~4 (Section~\ref{sec:lean}). The formal proofs were also
written by AI agents, but they are checked by the Lean kernel. The probabilistic part
(Lemmas~\ref{lem:anti} and~\ref{lem:expdef}), the lemma on random 3-CNF, the corollary and
Proposition~\ref{prop:dp} are not formalised.

\paragraph{No human expert review.} At the time of writing, \textbf{no human expert has
reviewed this work}. The author is a computer science student. Readers should treat the
results accordingly, and corrections are welcome at the address on the first page.

\paragraph{Sources and novelty.} The definitions and the question are quoted from the text
of~\cite{DL}. Among the works cited, the agents read \cite{DL} in full; the
preprints~\cite{OpenAI,Claude} and the exposition~\cite{ODonnell} are cited for context
only, and the present paper does not rely on their content beyond what~\cite{DL} reports.
The classical references were supplied by the AI agents from their knowledge of the
literature and were not re-read for this paper. The novelty assessment rests on searches made by AI
agents, not on a systematic survey.

\paragraph{Responsibility.} The author takes responsibility for the content of this paper,
including any error that remains.

\section{Conclusion and open problems}\label{sec:concl}

Dias and Lima asked for a proof that measurements cannot absorb the cut in clause-product
interpolation schemes. We proved a quantitative form of this statement: at some evaluation
point, the number of measurement rows plus the Taylor order exceeds
$c_b\sep(F)/\sqrt{s_t}$, or the groups are too large; and for variable-uniform deformations
the factor $\sqrt{s_t}$ disappears. On bounded-occurrence formulas with linear separators,
and on random 3-CNF of density above $3$, the design described in \cite[Section~9]{DL} does
not exist, and every scheme applied to the formula itself has product-type charge
$2^{\Omega(n^{2/3})}$. We also observed that the DP-type model cannot certify more than
$2^{\lceil\sqrt n\rceil}$ on any formula.

These are modest negative results about one abstraction of one route of one preprint. They
do not concern the preprint's parity form or tensor contraction, they are lower bounds on
the charges of a model and not on running time, and they have no bearing on ETH or on
P versus NP.

\paragraph{Open problems.}
\begin{enumerate}
\item \emph{Tightness of the $\sqrt s$ factor.} Is there a formula with
 $\sep(F)=\Omega(n)$ and a scheme with $\rho_t+q_t=O(n/\sqrt{s_t})$ at every point and
 $s_t=o(n)$? By Theorem~\ref{thm:uniform} its deformation must depend on the clause. The
 factor comes only from the fluctuations of a random witness map, and a different choice of
 witnesses, or a different minor, might remove it. If it can be removed, the product-type
 bound becomes $2^{\Omega(n)}$; if it cannot, the truth could be as low as
 $2^{\tilde\Theta(n^{2/3})}$.
 No upper bound better than the preprint's is known to us in the product-type model.
\item \emph{The preprint's actual inputs.} Prove $\sep(F^\star)=\Omega(n^\star)$, or its
 weighted analogue, for the bounded-occurrence formula $F^\star$ built with Tovey's
 cycles~\cite{Tov84}, and for the leaves of the sparsification. This would extend
 Corollary~\ref{cor:main} to the setting of \cite[Theorem~8.3]{DL}. We have not done it.
\item \emph{A meaningful DP-type model.} By Proposition~\ref{prop:dp} the DP-type charge
 must be repaired before the question can be asked in it. A natural candidate replaces
 $2^{w_t}$ by $2^{W_t}$, where $W_t$ is the minimum, over tree decompositions of the graph
 on measured groups used in~\cite{DL}, of the largest total number of rows in a bag. Is
 $\max(s_t,W_t,q_t)\ge n^{1/2+\delta}$ at some point, on formulas with linear separators?
 This is open. Partial results in this direction exist in the project repository; they are
 conditional and unreviewed, and we do not claim them here.
\item \emph{Formalisation.} Formalise Lemmas~\ref{lem:anti} and~\ref{lem:expdef}, so that
 Theorem~\ref{thm:general} is machine-checked without hypothesis, and remove the hypothesis
 $n\le M$ from the Lean statements.
\item \emph{Constants.} Improve $c_0$; the experiments suggest that
 Lemma~\ref{lem:anti} is far from sharp.
\item \emph{Beyond Definition~5.1.} Do similar limits hold for the parity form or the tensor
 contraction of the preprint? As explained in \cite[Section~5.1]{DL}, these break the link
 between the number of occurrences and the degree of exact recovery, on which
 Lemma~\ref{lem:hermite} and our comparison of degrees rest. Our method does not apply to
 them as it stands.
\end{enumerate}

\appendix
\section{Two properties of random 3-CNF (after Dias and Lima)}\label{app:random}

For completeness we reproduce the proofs of the two facts used in the proof of
Lemma~\ref{lem:random}. They are Lemmas~A.1 and~A.2 of~\cite{DL}, and the proofs follow
theirs. $F$ has $n$ variables and $m=\Delta n$ independent clauses, each on three distinct
uniformly random variables with random signs; $\log$ is the natural logarithm, and ``with
high probability'' means with probability tending to $1$ as $n\to\infty$.

\begin{lemma}[{\cite[Lemma~A.1]{DL}}]\label{lem:A1}
If $\Delta>3$, then with high probability, for every set $S$ of variables with
$n/3\le|S|\le2n/3$, at least $m/3$ clauses contain both a variable of $S$ and a variable
outside $S$.
\end{lemma}

\begin{proof}
Fix $S$ with $|S|=\alpha n$, $\alpha\in[1/3,2/3]$. The variable set of a clause is a uniform
$3$-subset of $[n]$, so it lies inside $S$ with probability
$\binom{|S|}{3}/\binom n3\le\alpha^3$, and inside the complement of $S$ with probability at
most $(1-\alpha)^3$. The function $\alpha^3+(1-\alpha)^3$ is convex, so on $[1/3,2/3]$ it is
at most its value $1/27+8/27=1/3$ at the endpoints. Hence each clause, independently,
contains a variable of $S$ and a variable outside $S$ with probability at least $2/3$, and
the number of such clauses stochastically dominates a binomial variable
$\mathrm{Bin}(m,2/3)$. By the Chernoff--Hoeffding bound~\cite[Theorem~1]{Hoe63},
\[
\PP\bigl(\mathrm{Bin}(m,2/3)\le m/3\bigr)\ \le\ \exp\bigl(-m\,\mathrm{KL}(1/3\,\|\,2/3)\bigr)
=2^{-m/3},
\]
because $\mathrm{KL}(1/3\,\|\,2/3)=\frac13\ln\frac12+\frac23\ln2=\frac13\ln2$. There are at
most $2^n$ sets $S$, so the probability that some $S$ fails is at most
$2^{n-\Delta n/3}$, which tends to $0$ because $\Delta>3$.
\end{proof}

\begin{lemma}[{\cite[Lemma~A.2]{DL}}]\label{lem:A2}
Let $C_\Delta=3e^2\Delta$. With high probability, for every $1\le j\le n$, every set of $j$
variables meets at most $C_\Delta\,j\log(en/j)$ clauses.
\end{lemma}

\begin{proof}
Fix a set $S$ of $j$ variables and let $L_S$ be the number of pairs (clause, variable of $S$
in that clause). The number of clauses that meet $S$ is at most $L_S$. For one clause, the
number of its variables in $S$ is hypergeometric: three draws without replacement from $n$
variables, $j$ of which are in $S$. By a theorem of Hoeffding~\cite[Theorem~4]{Hoe63}, its
moment generating function is at most that of $\mathrm{Bin}(3,j/n)$. The clauses are
independent, so for every $\lambda>0$,
\[
\E e^{\lambda L_S}\ \le\ \bigl(1+\tfrac jn(e^\lambda-1)\bigr)^{3m}\ \le\ \exp\bigl(\mu(e^\lambda-1)\bigr),
\qquad \mu=\frac{3mj}{n}=3\Delta j .
\]
By Markov's inequality with $\lambda=\ln(x/\mu)$, for every $x>\mu$,
\[
\PP(L_S\ge x)\ \le\ \exp\bigl(-x\ln(x/\mu)+x-\mu\bigr)\ \le\ \Bigl(\frac{e\mu}{x}\Bigr)^{x}.
\]
Take $x=C_\Delta\,j\log(en/j)$. Since $j\le n$, $\log(en/j)\ge1$, so $x\ge C_\Delta j>\mu$ and
$e\mu/x=1/(e\log(en/j))\le1/e$. Hence
$\PP(L_S\ge x)\le e^{-x}=(en/j)^{-C_\Delta j}$. There are $\binom nj\le(en/j)^j$ sets of size
$j$, so the probability that some set of size $j$ fails is at most
$(en/j)^{-(C_\Delta-1)j}\le(en/j)^{-2j}$. For $j\le\sqrt n$ this is at most $(e^2n)^{-j}$, and
for $j>\sqrt n$ it is at most $e^{-2j}$. Summing over $1\le j\le n$ gives at most
$2/(e^2n)+2e^{-2\sqrt n}$, which tends to $0$.
\end{proof}

\begin{thebibliography}{99}

\bibitem{AVW18} J. Alman, V. Vassilevska Williams. Limits on all known (and some unknown)
 approaches to matrix multiplication. \emph{FOCS} 2018.

\bibitem{AFL15} A. Ambainis, Y. Filmus, F. Le Gall. Fast matrix multiplication: limitations
 of the Coppersmith--Winograd method. \emph{STOC} 2015.

\bibitem{BSI10} E. Ben-Sasson, R. Impagliazzo. Random CNF's are hard for the polynomial
 calculus. \emph{Comput. Complexity} 19, 2010.

\bibitem{BSW01} E. Ben-Sasson, A. Wigderson. Short proofs are narrow---resolution made
 simple. \emph{J. ACM} 48(2), 2001.

\bibitem{Ber97} B. Berger. The fourth moment method. \emph{SIAM J. Comput.} 26(4), 1997.

\bibitem{Bod93} H. L. Bodlaender. A tourist guide through treewidth. \emph{Acta Cybernetica}
 11, 1993.

\bibitem{CS88} V. Chv\'atal, E. Szemer\'edi. Many hard examples for resolution.
 \emph{J. ACM} 35(4), 1988.

\bibitem{Claude} Anthropic (Claude). A $2^{o(n)}$ time algorithm for 3-SAT. Preprint,
 Hexagon 2610.00577v1, 26 October 2026.

\bibitem{DL} F. Dias, Y. Lima. Limits and constants of OpenAI's claimed ETH refutation.
 arXiv:2611.00517v1 [cs.CC], 31 October 2026.

\bibitem{DKSS13} Z. Dvir, S. Kopparty, S. Saraf, M. Sudan. Extensions to the method of
 multiplicities, with applications to Kakeya sets and mergers. \emph{SIAM J. Comput.} 42(6),
 2013.

\bibitem{Erd45} P. Erd\H{o}s. On a lemma of Littlewood and Offord. \emph{Bull. Amer. Math.
 Soc.} 51, 1945.

\bibitem{GLR82} I. Gohberg, P. Lancaster, L. Rodman. \emph{Matrix Polynomials}. Academic
 Press, 1982.

\bibitem{Hoe63} W. Hoeffding. Probability inequalities for sums of bounded random variables.
 \emph{J. Amer. Statist. Assoc.} 58(301), 1963.

\bibitem{IP01} R. Impagliazzo, R. Paturi. On the complexity of $k$-SAT. \emph{J. Comput.
 System Sci.} 62(2), 2001.

\bibitem{IPZ01} R. Impagliazzo, R. Paturi, F. Zane. Which problems have strongly exponential
 complexity? \emph{J. Comput. System Sci.} 63(4), 2001.

\bibitem{Kol58} A. Kolmogorov. Sur les propri\'et\'es des fonctions de concentrations de
 M.~P.~L\'evy. \emph{Ann. Inst. H. Poincar\'e} 16, 1958.

\bibitem{LT79} R. J. Lipton, R. E. Tarjan. A separator theorem for planar graphs.
 \emph{SIAM J. Appl. Math.} 36(2), 1979.

\bibitem{Lit30} J. E. Littlewood. On bounded bilinear forms in an infinite number of
 variables. \emph{Quart. J. Math.} 1, 1930.

\bibitem{LO43} J. E. Littlewood, A. C. Offord. On the number of real roots of a random
 algebraic equation (III). \emph{Rec. Math. [Mat. Sbornik] N.S.} 12(54), 1943.

\bibitem{Mathlib} The mathlib Community. The Lean mathematical library. \emph{CPP} 2020.

\bibitem{Lean4} L. de Moura, S. Ullrich. The Lean 4 theorem prover and programming language.
 \emph{CADE} 2021.

\bibitem{ODonnell} R. O'Donnell. Some exposition of ``The Exponential Time Hypothesis Is
 False''. Manuscript, 2026.
 \url{https://www.cs.cmu.edu/~odonnell/some_eth_false_exposition.pdf}

\bibitem{OpenAI} OpenAI. The exponential time hypothesis is false. Preprint, 17 October 2026,
 revised 25 October 2026.

\bibitem{RS86} N. Robertson, P. D. Seymour. Graph minors. II. Algorithmic aspects of
 tree-width. \emph{J. Algorithms} 7(3), 1986.

\bibitem{Rog61} B. A. Rogozin. On the increase of dispersion of sums of independent random
 variables. \emph{Theory Probab. Appl.} 6(1), 1961.

\bibitem{Tov84} C. A. Tovey. A simplified NP-complete satisfiability problem.
 \emph{Discrete Appl. Math.} 8, 1984.

\end{thebibliography}

\end{document}
